\chapter{METODE PENELITIAN}

\section{Desain Penelitian}
Penelitian ini dilakukan dengan mengadopsi metode \textit{Design and Development} (D\&D) yang dirumuskan oleh \textcite{ellisGuideNoviceResearchers2010}. Metode D\&D didefinisikan sebagai pendekatan sistematis untuk merancang, mengembangkan, dan mengevaluasi suatu produk atau sistem (artefak) sebagai solusi atas permasalahan yang spesifik. Berbeda dengan metode penelitian tradisional yang umumnya berfokus pada pengujian hipotesis atau observasi fenomena, penelitian D\&D menempatkan pengembangan artefak teknologi itu sendiri sebagai kontribusi utama terhadap khazanah pengetahuan. Tujuan utamanya adalah untuk menjembatani kesenjangan antara landasan teoritis dan implementasi praktis melalui penciptaan solusi yang dapat diuji validitas dan efektivitasnya.

Metode ini dipilih karena kesesuaiannya dengan karakteristik penelitian ini yang berfokus pada pengembangan prototipe sistem agentik. Menurut \textcite{ellisGuideNoviceResearchers2010}, pendekatan D\&D menyediakan kerangka kerja yang terstruktur untuk memandu proses teknis mulai dari konseptualisasi hingga validasi sistem. Hal ini dinilai relevan mengingat luaran penelitian masih terbatas dalam bentuk prototipe yang berhenti pada tahap validasi fungsional dan kinerja, sehingga diperlukan metodologi yang menekankan pada iterasi desain dan pembuktian konsep (\textit{proof of concept}). Tahapan pelaksanaan penelitian mengacu pada enam fase dalam model D\&D, yaitu identifikasi masalah, deskripsi tujuan, perancangan dan pengembangan, pengujian, evaluasi, dan penyampaian hasil. Alur tahapan tersebut divisualisasikan pada Gambar \ref{fig:dnd}.

\begin{figure}[H] 
    \centering
    \includegraphics[width=1\textwidth]{bab3/dnd.png}
    \caption[Tahapan penelitian metode \textit{Design and Development} (D\&D)]{Tahapan penelitian metode \textit{Design and Development} (D\&D) \parencite{ellisGuideNoviceResearchers2010}}
    \label{fig:dnd}
\end{figure}

\section{Identifikasi Masalah}


% \begin{table}[H]
%     \centering
%     \caption{Perangkat Pendukung Penelitian}
%     \label{tab:perangkat_pendukung}
%     \small
%     \setstretch{1.0}
%     \begin{tabular}{|c|l|p{9cm}|}
%     \hline
%     \multicolumn{3}{|c|}{\textbf{Spesifikasi Laptop}} \\ \hline
%     \textbf{No.} & \textbf{Komponen} & \textbf{Spesifikasi} \\ \hline
%     1. & Prosesor & AMD Ryzen™ 7 7840HS \\ \hline
%     2. & Memori & 16 GB DDR5-5600 \\ \hline
%     3. & Grafik & NVIDIA® GeForce RTX™ 4060 8GB \\ \hline
%     \multicolumn{3}{|c|}{\textbf{Perangkat Lunak}} \\ \hline
%     \textbf{No.} & \textbf{Nama} & \textbf{Fungsi} \\ \hline
%     1. & WSL2 & Subsistem Linux untuk Windows sebagai lingkungan pengembangan sistem agentik. \\ \hline
%     2. & Docker & Platform kontainerisasi aplikasi. \\ \hline
%     3. & uv & Manajer library dan virtual environtment untuk Python. \\ \hline
%     4. & Next.js & Framework aplikasi web. \\ \hline
%     5. & LangChain & Framework pengembangan aplikasi LLM. \\ \hline
%     6. & LangGraph & Framework untuk pengembangan sistem agentik. \\ \hline
%     7. & LangSmith & Platform monitoring dan evaluasi LLM. \\ \hline
%     8. & AgentEval & Library evaluasi sistem agentik dari LangSmith. \\ \hline
%     9. & DeepEval & Library penyedia metrik evaluasi LLM-as-a-judge untuk sistem agentik. \\ \hline
%     10. & RAGAS & Library penyedia metrik evaluasi LLM-as-a-judge untuk RAG. \\ \hline
%     11. & Qdrant & Basis data vektor untuk pencarian semantik. \\ \hline
%     12. & Google Gen AI & Layanan penyedia API model bahasa Gemini. \\ \hline
%     13. & Voyage AI & Layanan penyedia API model embedding dan reranker. \\ \hline
%     14. & Kaggle Kernel & Platform komputasi untuk pelatihan model. \\ \hline
%     15. & Supabase & Layanan penyedia penyimpanan data. \\ \hline
%     16. & Cloudflare R2 & Layanan penyedia penyimpanan objek. \\ \hline
%     17. & FastEmbed & Library inferensi model embedding untuk RAG. \\ \hline
%     18. & LlamaIndex & Framework untuk pengembangan RAG. \\ \hline
%     19. & Playwright & Framework untuk automasi web untuk \textit{crawling}. \\ \hline
%     20. & ScrapeGraphAI & Library untuk ekstraksi data web berbasis LLM. \\ \hline
%     21. & Google Cloud Run & Layanan deployment cloud container . \\ \hline
%     22. & Vercel & Layanan hosting aplikasi web. \\ \hline
%     \end{tabular}
% \end{table}

\section{Deskripsi Tujuan}
Berdasarkan identifikasi masalah yang telah diuraikan, tujuan utama dari penelitian ini adalah mengembangkan sistem agentik multimodal yang mampu mengidentifikasi penyakit tanaman serta memberikan penjelasan diagnosis yang berbasis bukti. Sistem ini dirancang sebagai solusi untuk mengatasi keterbatasan metode identifikasi konvensional yang kurang adaptif. Secara rinci, tujuan pengembangan dalam penelitian ini dijabarkan sebagai berikut:
\begin{enumerate}[nosep]
    \item Mengembangkan arsitektur sistem agentik menggunakan pola \textit{Reasoning and Acting} (ReAct) yang memungkinkan model bahasa besar untuk memilih dan menggunakan alat bantu secara adaptif dalam proses diagnosis.
    \item Mengintegrasikan model deteksi objek YOLOv11 dan OWLv2 untuk melokalisasi gejala penyakit pada bagian tanaman secara spesifik.
    \item Mengimplementasikan modul pencarian multimodal berbasis model SCOLD untuk memungkinkan identifikasi penyakit melalui pencarian kemiripan visual dan tekstual pada basis data vektor.
    \item Menerapkan mekanisme \textit{Retrieval-Augmented Generation} (RAG) untuk memperkaya respon sistem dengan informasi manajemen penyakit yang faktual dari sumber pengetahuan eksternal.
    \item Merancang antarmuka pengguna berbasis web yang memfasilitasi interaksi dua arah antara pengguna dan sistem melalui teks dan gambar.
    \item Mengevaluasi kinerja sistem agentik yang dikembangkan menggunakan pendekatan \textit{LLM-as-a-judge} untuk mengukur akurasi diagnosis dan relevansi saran yang diberikan.
\end{enumerate}

\section{Perancangan dan Pengembangan}
Tahap perancangan dan pengembangan difokuskan pada realisasi teknis dari sistem agentik yang diusulkan. Pada bagian ini, arsitektur sistem dan implementasi komponen-komponen utamanya dijabarkan secara rinci. Pembahasan diawali dengan perancangan sistem agentik yang mengadopsi pola ReAct dan pengembangan agen LLM yang berfungsi sebagai pengendali logika sistem. Selanjutnya, uraian dilanjutkan pada pengembangan modul-modul spesifik yang mendukung kapabilitas sistem, meliputi modul deteksi objek tanaman, modul pencarian multimodal gambar-teks, serta modul pencarian informasi. Akhirnya, pengembangan antarmuka pengguna dipaparkan sebagai komponen interaksi antara pengguna dan sistem yang telah dibangun.

\subsection{Sistem Agentik}
Sistem agentik dalam penelitian ini dirancang dengan mengadopsi pola \textit{Reasoning and Acting} (ReAct). Implementasi sistem dilakukan menggunakan framework LangChain melalui library LangGraph \parencite{langchainaiLangChain2025}. Arsitektur sistem yang dikembangkan bersifat multimodal, di mana agen mampu menerima dan memproses input berupa citra tanaman serta instruksi tekstual dari pengguna secara bersamaan. Gambar \ref{fig:arsitektur_sistem_agentik} menyajikan arsitektur sistem agentik beserta spesifikasi input dan output pada setiap alat yang terintegrasi.

\begin{figure}[H]
    \centering
    \includegraphics[width=1\textwidth]{bab3/arsitektursistemagentik.png}
    \caption{Arsitektur sistem agentik dengan kemampuan multimodal}
    \label{fig:arsitektur_sistem_agentik}
\end{figure}

Dalam implementasi teknis menggunakan LangGraph, konsep teoritis ReAct yang diajukan oleh \textcite{yaoReActSynergizingReasoning2023} diterjemahkan ke dalam struktur graf berstatus (\textit{stateful graph}). Pada desain aslinya, pola ReAct diimplementasikan secara linier melalui teknik \textit{prompting}, di mana LLM menghasilkan teks terstruktur berisi penalaran (\textit{thought}) dan tindakan (\textit{action}) yang kemudian diurai oleh lingkungan eksternal. Di lain sisi, LangGraph memodelkan agen sebagai sistem yang terdiri dari \textit{node} (simpul) dan \textit{edge} (sisi). 

Struktur ini memetakan siklus Thought-Action-Observation ke dalam dua komponen utama, yaitu \textit{node} Model yang berisi agen dan \textit{node} Tools yang berisi alat. \textit{node} Model bertanggung jawab untuk menentukan apakah suatu tugas memerlukan penggunaan alat berdasarkan konteks yang diberikan. Jika alat diperlukan, LLM akan menghasilkan permintaan pemanggilan alat (\textit{tool call}) yang valid secara skema untuk dieksekusi secara otomatis oleh \textit{node} Tools. Hasil eksekusi tersebut dikembalikan ke \textit{node} Model sebagai observasi untuk dianalisis lebih lanjut hingga solusi akhir tercapai. 

Status (\textit{state}) agen dikelola sebagai objek yang menyimpan riwayat pesan, yang berfungsi sebagai memori jangka pendek atau \textit{scratchpad}. Alur kerja dimulai dari \textit{node} Model, di mana LLM memproses input dan menentukan langkah selanjutnya. Jika model memutuskan untuk menggunakan alat, alur dialihkan melalui \textit{conditional edge} menuju \textit{node} Tools. Setelah alat dieksekusi, hasilnya dicatat ke dalam state sebagai observasi, dan kendali dikembalikan ke \textit{node} Model untuk mengevaluasi informasi baru tersebut. Proses iteratif ini berlanjut hingga model menilai bahwa informasi yang dikumpulkan telah memadai untuk memberikan jawaban akhir (END). Visualisasi perbandingan antara desain teoritis ReAct dengan implementasi berbasis graf pada LangGraph disajikan pada Gambar \ref{fig:react_langgraph_comparison}.

\begin{figure}[H]
    \centering
    \includegraphics[width=1\textwidth]{bab3/reactlanggraph.png}
    \caption[Perbandingan antara (a) desain teoritis ReAct dan (b) implementasi berbasis graf pada LangGraph]{Perbandingan antara (a) desain teoritis ReAct oleh \textcite{yaoReActSynergizingReasoning2023} dan (b) implementasi berbasis graf pada LangGraph}
    \label{fig:react_langgraph_comparison}
\end{figure}

\subsection{Agen LLM}
Bagian ini menjelaskan spesifikasi dan konfigurasi agen LLM yang dikembangkan untuk mengimplementasikan pola ReAct pada sistem identifikasi penyakit tanaman.

\subsubsection{Konfigurasi}
Agen dalam penelitian ini digerakkan oleh Gemini 3 Flash sebagai LLM utama. Pemilihan model ini didasarkan pada karakteristik model yang menawarkan kecepatan inferensi tinggi dan kemampuan penalaran (\textit{reasoning}) serta penggunaan alat (\textit{tool-use}) yang sangat baik untuk tugas-tugas agentik. Konfigurasi model diatur dengan nilai \textit{temperature} 0 guna memastikan luaran yang bersifat deterministik dan konsisten, sehingga meminimalkan variasi jawaban yang tidak relevan dalam konteks teknis identifikasi penyakit.

Melalui implementasi arsitektur LangGraph, agen memiliki memori jangka pendek (\textit{short-term memory}) yang dikelola melalui objek \textit{State} (status) yang bersifat \textit{stateful}. Status ini berfungsi sebagai buku catatan (\textit{scratchpad}) yang menyimpan seluruh riwayat interaksi dan informasi yang dikumpulkan selama satu sesi tugas berlangsung. Dengan menyimpan status secara terstruktur, agen memiliki konteks terhadap langkah-langkah yang telah diambil sebelumnya, misalnya hasil deteksi objek atau temuan dari pencarian internet, yang kemudian digunakan sebagai dasar untuk menentukan langkah tindakan selanjutnya hingga solusi akhir tercapai. Tabel \ref{tab:state_components} menyajikan rincian variabel status yang dikelola dalam sistem.

\begin{table}[H]
    \centering
    \caption{Komponen Status (\textit{State}) pada Agen LLM}
    \label{tab:state_components}
    \small
    \setstretch{1.0}
    \begin{tabular}{|l|l|p{4cm}|}
    \hline
    \multicolumn{1}{|c|}{\textbf{Variabel Status}} & \multicolumn{1}{c|}{\textbf{Tipe Data}} & \multicolumn{1}{c|}{\textbf{Deskripsi}} \\ \hline
    \textit{messages} & \textit{list[AnyMessage]} & Menyimpan riwayat seluruh pesan interaksi (percakapan) antara pengguna dan agen sebagai memori jangka pendek (\textit{scratchpad}). \\ \hline
    \textit{images} & \textit{list[dict]} & Menyimpan daftar gambar yang diunggah beserta metadata terkait untuk mendukung penanganan beberapa citra sekaligus. \\ \hline
    \textit{detections} & \textit{list[dict]} & Menyimpan data hasil deteksi objek yang mencakup label, skor kepercayaan, dan koordinat kotak pembatas (\textit{bounding box}). \\ \hline
    \textit{plant\_disease\_classifications} & \textit{list[dict]} & Menyimpan hasil identifikasi penyakit tanaman yang mencakup label prediksi, tingkat kepercayaan, dan metadata pendukung lainnya. \\ \hline
    \textit{visualization\_urls} & \textit{list[str]} & Menyimpan daftar tautan publik dari citra visualisasi hasil deteksi yang telah diunggah ke penyimpanan Cloudflare R2. \\ \hline
    \textit{current\_image\_url} & \textit{str} & Menyimpan URL dari gambar yang sedang menjadi fokus analisis aktif dalam proses iterasi berjalan. \\ \hline
    \textit{jump\_to} & \textit{JumpTo} & Atribut kontrol internal untuk mengelola navigasi alur kerja (\textit{routing}) antar simpul (\textit{node}) dalam arsitektur graf LangGraph. \\ \hline
    \end{tabular}
\end{table}

Untuk memastikan stabilitas, keamanan, dan fungsionalitas yang optimal, agen dilengkapi dengan serangkaian \textit{middleware}. \textit{Middleware} dalam arsitektur ini berfungsi sebagai lapisan perantara yang membungkus pemanggilan model maupun alat (\textit{tools}) untuk pemrosesan tambahan secara otomatis sebelum dan sesudah eksekusi dilakukan. Tabel \ref{tab:middleware_components} merinci komponen \textit{middleware} yang diintegrasikan ke dalam sistem.

\begin{table}[H]
    \centering
    \caption{Komponen \textit{Middleware} pada Agen LLM}
    \label{tab:middleware_components}
    \small
    \setstretch{1.0}
    \begin{tabular}{|l|p{7cm}|}
    \hline
    \multicolumn{1}{|c|}{\textbf{Middleware}} & \multicolumn{1}{c|}{\textbf{Deskripsi Fungsi}} \\ \hline
    \textit{get\_system\_prompt} & Menyediakan instruksi dasar, batasan peran, dan panduan penggunaan alat kepada agen. \\ \hline
    \textit{handle\_tool\_errors} & Menangkap kegagalan eksekusi pada \textit{tools} dan mengembalikan pesan kesalahan yang informatif agar agen dapat melakukan koreksi secara mandiri. \\ \hline
    \textit{model\_retry\_middleware} & Melakukan percobaan ulang otomatis jika terjadi kegagalan koneksi atau batasan kuota (\textit{rate limit}) pada API model. \\ \hline
    \textit{tool\_retry\_middleware} & Menangani percobaan ulang jika eksekusi alat eksternal mengalami gangguan teknis sementara. \\ \hline
    \textit{model\_fallback\_middleware} & Mekanisme cadangan yang secara otomatis mengalihkan tugas ke model yang lebih tinggi, yakni Gemini 3 Pro apabila model utama gagal merespons. \\ \hline
    \end{tabular}
\end{table}

\subsubsection{\textit{Prompt Engineering}}
Desain instruksi (\textit{prompt engineering}) pada agen LLM dirancang dengan mengadopsi kerangka kerja CO-STAR (Context, Objective, Style, Tone, Audience, Response). Kerangka kerja ini digunakan untuk memastikan bahwa model mendapatkan konteks yang lengkap dan terstruktur guna menghasilkan respons yang akurat, relevan, dan konsisten \parencite{govtechsingaporeCOSTARPromptEngineering2024}. Gambar \ref{fig:costar_framework} menyajikan visualisasi dari kerangka kerja CO-STAR.

\begin{figure}[H]
    \centering
    \includegraphics[width=1\textwidth]{bab3/costar.png}
    \caption{Visualisasi kerangka kerja CO-STAR}
    \label{fig:costar_framework}
\end{figure}

CO-STAR merupakan akronim dari enam komponen utama yang dijelaskan sebagai berikut:
\begin{enumerate}[nosep]
    \item \textit{Context} (Konteks): Memberikan informasi latar belakang mengenai peran agen sebagai asisten identifikasi penyakit tanaman, batasan domain penelitian, serta kapabilitas sistem.
    \item \textit{Objective} (Tujuan): Menetapkan tujuan utama agen untuk menghasilkan identifikasi yang berbasis bukti, jujur terhadap ketidakpastian, dan memberikan panduan yang berlandaskan data.
    \item \textit{Style} (Gaya): Menentukan cara penulisan respons agar mengalir secara alami, eksplisit dalam penyebutan sumber bukti, serta ringkas tanpa pengulangan informasi.
    \item \textit{Tone} (Nada): Mengatur nada bicara agen agar bersifat praktis, tenang, dan profesional tanpa memberikan peringatan yang berlebihan.
    \item \textit{Audience} (Audiens): Mengidentifikasi target pengguna sistem yang mencakup petani, penghobi tanaman, hingga tenaga profesional di bidang pertanian dengan tingkat pemahaman teknis yang bervariasi.
    \item \textit{Response} (Respons): Menentukan struktur luaran yang diharapkan, mulai dari observasi visual, integrasi konteks pengguna, hingga hipotesis diagnosis dan saran tindak lanjut.
\end{enumerate}

Selain komponen standar CO-STAR, desain prompt pada sistem ini juga dilengkapi dengan bagian penjelasan alat (\textit{Tools Reference}). Bagian ini berfungsi sebagai dokumentasi teknis yang menjelaskan fungsionalitas, parameter, dan skenario penggunaan dari setiap alat yang terintegrasi, seperti pencarian web, pencarian basis pengetahuan, dan model deteksi visual. Dengan adanya penjelasan alat yang komprehensif, agen dapat memahami cara berinteraksi dengan lingkungan eksternal secara tepat.

\subsection{Modul Deteksi Objek Tanaman}
Modul deteksi objek tanaman dirancang sebagai tahap awal dalam alur kerja identifikasi penyakit untuk melakukan lokalisasi terhadap bagian tanaman yang relevan, seperti daun, buah, atau batang. Modul ini memungkinkan agen untuk memperoleh kotak pembatas (\textit{bounding box}) dari objek yang terdeteksi, yang kemudian digunakan oleh modul pencarian gambar–teks untuk melakukan pencarian berbasis wilayah (\textit{region-based search}). Pendekatan ini bertujuan untuk mereduksi derau latar belakang (\textit{background noise}) serta memitigasi gangguan dari objek non-relevan lainnya dalam satu citra, sehingga meningkatkan akurasi pengambilan informasi.  Alur kerja dari modul deteksi objek tanaman divisualisasikan pada Gambar \ref{fig:plant_detection_workflow}.

Implementasi modul ini menggunakan pendekatan hibrida yang menggabungkan dua alat dengan karakteristik yang berbeda, yakni alat deteksi \textit{closed-set} dan alat deteksi \textit{open-vocabulary}. Alat deteksi \textit{closed-set} menggunakan arsitektur YOLOv11 dari \textcite{khanamYOLOv11OverviewKey2024} yang dioptimalkan khusus untuk deteksi daun, mengingat daun merupakan bagian tanaman yang paling dominan dalam memperlihatkan gejala penyakit. Sementara itu, alat deteksi \textit{open-vocabulary} menggunakan model OWLv2 dari \textcite{mindererScalingOpenVocabularyObject2024} untuk mendeteksi bagian tanaman lain melalui perintah teks (\textit{text prompt}) secara fleksibel. YOLOv11 menawarkan kecepatan inferensi dan akurasi yang sangat tinggi, namun untuk kelas terbatas yang spesifik berupa daun, sedangkan OWLv2 memberikan fleksibilitas deteksi objek kelas selain daun, seperti buah atau batang, namun dengan beban komputasi yang lebih besar.

\begin{figure}[H]
    \centering
    \includegraphics[width=1\textwidth]{bab3/moduldeteksitanaman.png}
    \caption{Alur Kerja Modul Deteksi Objek Tanaman}
    \label{fig:plant_detection_workflow}
\end{figure}

\subsubsection{Alat Deteksi \textit{Closed-Set}}
Alat deteksi \textit{closed-set} dalam modul ini diimplementasikan menggunakan arsitektur YOLOv11. Pemilihan YOLOv11 didasarkan pada performanya yang unggul dalam menyeimbangkan antara kecepatan inferensi dan akurasi deteksi dibandingkan versi sebelumnya. Model ini dilatih secara khusus untuk mengenali objek daun di lapangan.

\paragraph{Persiapan Dataset Pemodelan}
Dataset yang digunakan untuk melatih model deteksi daun diadopsi dari dataset PlantDoc dari \textcite{singhPlantDocDatasetVisual2020}. Dataset ini dipilih karena mencakup variasi spesies tanaman dan kondisi penyakit yang luas dalam lingkungan alami. Secara orisinal, dataset ini terdiri dari 2.569 citra yang merepresentasikan 13 spesies tanaman dan terbagi ke dalam 30 kelas kondisi. Dalam penelitian ini, dilakukan prapemrosesan berupa penyatuan seluruh anotasi kelas yang ada menjadi satu kelas tunggal, yakni "\textit{leaf}" (daun). Hal ini dilakukan untuk memfokuskan model pada tugas lokalisasi daun sebagai \textit{Region of Interest} (RoI) utama. Distribusi data awal dari dataset PlantDoc sebelum penyatuan kelas disajikan pada Gambar \ref{fig:plantdoc_distribution} untuk memberikan gambaran keragaman data yang digunakan.

\begin{figure}[H]
    \centering
    \begin{subfigure}[b]{1\textwidth}
        \centering
        \includegraphics[width=\textwidth]{bab3/sampelplantdoc.png}
        \caption{}
        \label{fig:plantdoc_sample}
    \end{subfigure}
    
    \vspace{0.5cm}
    
    \begin{subfigure}[b]{1\textwidth}
        \centering
        \includegraphics[width=\textwidth]{bab3/plantdoc_distribution.png}
        \caption{}
        \label{fig:plantdoc_distribution_sub}
    \end{subfigure}
    \caption[Dataset PlantDoc: (a) Sampel citra dan (b) Distribusi kelas]{Dataset PlantDoc dari \textcite{singhPlantDocDatasetVisual2020}: (a) Sampel citra dan (b) Distribusi kelas}
    \label{fig:plantdoc_distribution}
\end{figure}

\paragraph{Pemodelan}
Pemodelan dilakukan dengan mengevaluasi tiga variasi skala model YOLOv11, yaitu YOLOv11n (nano), YOLOv11s (small), dan YOLOv11m (medium). Ketiga variasi model ini dipilih karena kesesuaiannya dengan target lingkungan inferensi berbasis CPU serta kebutuhan sumber daya komputasi pelatihan yang masih dapat diakomodasi oleh perangkat keras yang tersedia dalam penelitian ini. Ketiga varian ini dibandingkan untuk menentukan model yang paling optimal untuk diintegrasikan ke dalam sistem agentik. Seluruh proses pelatihan dilakukan menggunakan Ultralytics dengan konfigurasi hiperparameter bawaan, kecuali yang dirinci pada Tabel \ref{tab:yolo_hyperparameters}.

\begin{table}[H]
\centering
    \caption{Konfigurasi Hiperparameter Pelatihan YOLOv11}
    \label{tab:yolo_hyperparameters}
    \centering
    \setstretch{1.0}
    \begin{tabular}{|l|l|p{6cm}|}
\hline
    \multicolumn{1}{|c|}{\textbf{Parameter}} & \multicolumn{1}{c|}{\textbf{Nilai}} & \multicolumn{1}{c|}{\textbf{Keterangan}} \\ \hline
    Epoch & 240 & Jumlah iterasi pelatihan penuh. \\ \hline
    Input Size & 416 x 416 & Resolusi citra masukan sesuai anotasi PlantDoc. \\ \hline
    Random Seed & 42 & Menjamin keterulangan hasil eksperimen. \\ \hline
    Batch Size & 16 & Jumlah sampel per sub-iterasi. \\ \hline
    Patience & 20 & Batas iterasi tanpa peningkatan untuk \textit{early stopping}. \\ \hline
    \end{tabular}
\end{table}

Setelah proses pelatihan selesai, model diekspor ke format ONNX (\textit{Open Neural Network Exchange}) untuk meningkatkan portabilitas dan efisiensi saat dijalankan di lingkungan produksi. Proses ekspor dilakukan dengan pengaturan resolusi sesuai ukuran masukan (\textit{input size}) pada perangkat CPU guna memastikan kompatibilitas dengan keterbatasan sumber daya pada lingkungan server.

\subsubsection{Alat Deteksi Open-Vocabulary}
Alat deteksi \textit{open-vocabulary} menggunakan model OWLv2 \parencite{huggingfaceOWLv2} yang dikembangkan oleh Google \parencite{mindererScalingOpenVocabularyObject2024}. Sistem menggunakan varian model \textit{google/\allowbreak owlv2-base-\allowbreak patch16} yang telah dikonversi ke format ONNX. Model ini mengintegrasikan \textit{vision transformer} (ViT) dengan \textit{text encoder} untuk melakukan deteksi objek berdasarkan deskripsi tekstual secara \textit{zero-shot}. Gambar \ref{fig:arsitektur_owlv2} menyajikan visualisasi arsitektur dari model OWLv2.

\begin{figure}[H]
    \centering
    \includegraphics[width=1\textwidth]{bab3/arsitekturalatowlv2.png}
    \caption{Arsitektur Model OWLv2 yang Digunakan}
    \label{fig:arsitektur_owlv2}
\end{figure}

\subsection{Modul Pencarian Gambar-Teks Penyakit}
Modul pencarian gambar-teks penyakit dirancang sebagai komponen utama sistem untuk mengidentifikasi gejala penyakit tanaman melalui pendekatan pengambilan informasi multimodal. Implementasi modul ini menggunakan teknologi Qdrant yang mendukung penyimpanan dan pencarian vektor bernama jamak (\textit{multi-named vectors}) secara efisien. Melalui fitur ini, dataset yang berisi pasangan gambar dan teks dapat disematkan (\textit{embedded}) ke dalam satu koleksi Qdrant, di mana satu poin data menampung dua vektor sekaligus, yaitu vektor gambar dan vektor teks. Representasi fitur diekstraksi menggunakan model \textit{foundation} SCOLD \parencite{huggingfaceSCOLD} yang dikembangkan oleh \textcite{quocVisionLanguageFoundationModel2025} dengan arsitektur RoBERTa dari \textcite{liuRoBERTaRobustlyOptimized2019} sebagai \textit{encoder} teks dan arsitektur Swin-T dari \textcite{liuSwinTransformerHierarchical2021} sebagai \textit{encoder} gambar. Rincian lebih lanjut mengenai skema ini akan dibahas pada subbab berikutnya. Alur kerja dari modul pencarian gambar-teks divisualisasikan pada Gambar \ref{fig:image_text_search_workflow}.

\begin{figure}[H]
    \centering
    \includegraphics[width=1\textwidth]{bab3/modulgambarteks.png}
    \caption{Alur Kerja Modul Pencarian Gambar-Teks Penyakit Tanaman}
    \label{fig:image_text_search_workflow}
\end{figure}


\subsubsection{Persiapan Data Indeks}
Dataset PlantWild yang disusun oleh \textcite{weiBenchmarkingInthewildMultimodal2024} digunakan sebagai basis indeks pencarian pada modul ini. Dataset ini mencakup total 18.542 citra yang terdistribusi ke dalam 89 kelas kondisi kesehatan dan penyakit pada 13 spesies tanaman utama. Citra dalam dataset ini mencakup berbagai organ tanaman, termasuk daun, buah, dan batang. Klasifikasi tersebut terdiri dari 56 kelas penyakit dan 33 kelas kondisi sehat. Dalam penelitian ini, dilakukan pengambilan sampel secara acak sebanyak 3 citra dari setiap kelas untuk kebutuhan pengujian, sehingga diperoleh total 267 citra uji. Sebanyak 18.275 citra sisanya kemudian digunakan untuk diindeks ke dalam basis data vektor. Selain data citra, setiap kelas dalam dataset ini juga dilengkapi dengan 10 \textit{caption} yang mendeskripsikan karakteristik visual dari kelas tersebut. Gambar \ref{fig:plantwild_distribution} menyajikan distribusi kelas pada dataset PlantWild.
 
\begin{figure}[H]
    \centering
    \includegraphics[width=1\textwidth]{bab3/distribusiplantwild.png}
    \caption[Distribusi Kelas pada Dataset PlantWild]{Distribusi Kelas pada Dataset PlantWild dari \protect\textcite{weiBenchmarkingInthewildMultimodal2024}}
    \label{fig:plantwild_distribution}
\end{figure}

\subsubsection{\textit{Indexing}}
Dalam Qdrant, data dikelola dalam unit entitas yang disebut sebagai Poin (\textit{Points}). Setiap Poin terdiri dari tiga komponen utama, yaitu ID unik, vektor yang dapat berupa vektor tunggal atau jamak, dan \textit{payload} opsional untuk menyimpan metadata. Pada sistem ini, setiap entitas penyakit tanaman direpresentasikan sebagai Poin yang memiliki dua vektor padat (\textit{dense vectors}) bernama, yaitu vektor teks dan vektor citra. Citra referensi asli disimpan secara terpisah pada layanan penyimpanan objek Cloudflare R2 untuk efisiensi, sementara URL publiknya disimpan dalam \textit{payload}.

Tabel \ref{tab:qdrant_schema} menyajikan rincian skema \textit{Points} dan struktur \textit{payload} yang diimplementasikan. Skema ini dirancang untuk mendukung pencarian multimodal sekaligus penyaringan metadata yang presisi.

%TODO: perlu diperbagus tabelnya
\begin{table}[H]
    \centering
    \caption{Skema Poin (Vektor dan Payload) pada Indeks Qdrant}
    \label{tab:qdrant_schema}
    \small
    \setstretch{1.0}
    \resizebox{\textwidth}{!}{%
    \begin{tabular}{|l|l|l|p{7cm}|}
    \hline
    \multicolumn{1}{|c|}{\textbf{Komponen}} & \multicolumn{1}{c|}{\textbf{Kunci}} & \multicolumn{1}{c|}{\textbf{Tipe}} & \multicolumn{1}{c|}{\textbf{Deskripsi}} \\ \hline
    \multirow{2}{*}{\textbf{Vektor}} & \textit{text} & 512 & Representasi semantik deskripsi gejala (RoBERTa). \\ \cline{2-4} 
     & \textit{image} & 512 & Representasi fitur visual citra tanaman (Swin-T/SCOLD). \\ \hline
    \multirow{6}{*}{\textbf{Payload}} & \textit{plant\_name} & String & Nama spesies tanaman. Diindeks secara \textit{full-text} untuk filtrasi. \\ \cline{2-4} 
     & \textit{label} & String & Nama spesifik penyakit. Diindeks secara \textit{full-text}. \\ \cline{2-4} 
     & \textit{caption} & String & Deskripsi tekstual lengkap mengenai gejala. \\ \cline{2-4} 
     & \textit{image\_url} & String & Tautan publik citra di Cloudflare R2. \\ \cline{2-4} 
     & \textit{source\_id} & Integer & ID referensi internal. \\ \cline{2-4} 
     & \textit{filename} & String & Nama fail asli citra. \\ \hline
    \end{tabular}%
    }
\end{table}

Fitur penting dalam implementasi ini adalah penerapan indeks teks penuh (\textit{full-text index}) yang disediakan oleh Qdrant pada kolom \textit{payload} \textit{plant\_name} dan \textit{label}. Berbeda dengan pencocokan kata kunci eksak (\textit{keyword match}), indeks teks penuh memungkinkan agen untuk melakukan penyaringan (\textit{filtering}) yang lebih fleksibel, seperti pencocokan berdasarkan token kata (\textit{token-based matching}) dan ketidakpekaan terhadap huruf besar-kecil (\textit{case-insensitive}). Fitur ini secara signifikan meningkatkan efisiensi pencarian agen dengan membatasi ruang pencarian vektor hanya pada subset data yang relevan dengan konteks tanaman yang sedang dianalisis. Sebagai ilustrasi, ketika agen mendeteksi bahwa tanaman yang dianalisis adalah tomat, agen dapat menerapkan filter pada \textit{plant\_name} untuk hanya mencari di antara data penyakit tomat, sehingga hasil yang diperoleh lebih akurat dan relevan. Visualisasi mekanisme ini disajikan pada Gambar \ref{fig:payload_filtering}.

%TODO: ini gambar masih placeholder, sesuain konteksnya ke tanaman!
\begin{figure}[H]
    \centering
    \includegraphics[width=1\textwidth]{bab3/fulltextindexpayloadfield.png}
    \caption{Visualisasi Mekanisme Penyaringan}
    \label{fig:payload_filtering}
\end{figure}

\subsubsection{Strategi Pencarian}
Implementasi modul ini mendukung empat strategi pencarian multimodal yang memanfaatkan kemampuan representasi vektor ganda (teks dan citra) pada indeks Qdrant. Pemilihan strategi dilakukan secara dinamis oleh agen berdasarkan ketersediaan input dari pengguna (teks gejala atau citra tanaman) serta tujuan pencarian. Fleksibilitas ini dimungkinkan oleh arsitektur model SCOLD yang menyelaraskan (\textit{align}) representasi fitur teks dan citra ke dalam ruang laten yang sama, memungkinkan operasi pencarian lintas modal (\textit{cross-modal}). Keempat strategi tersebut diilustrasikan pada Gambar \ref{fig:multimodal_search_strategies} dan dijelaskan sebagai berikut:

\begin{figure}[H]
    \centering
    \begin{subfigure}[b]{0.45\textwidth}
        \centering
        \includegraphics[width=\textwidth]{bab3/strategipencarianteksteks.png}
        \caption{}
        \label{fig:strat_t2t}
    \end{subfigure}
    \hfill
    \begin{subfigure}[b]{0.45\textwidth}
        \centering
        \includegraphics[width=\textwidth]{bab3/strategipencariantekksgambar.png}
        \caption{}
        \label{fig:strat_t2i}
    \end{subfigure}
    
    \vspace{0.5cm}
    
    \begin{subfigure}[b]{0.45\textwidth}
        \centering
        \includegraphics[width=\textwidth]{bab3/strategipencariangambargambar.png}
        \caption{}
        \label{fig:strat_i2i}
    \end{subfigure}
    \hfill
    \begin{subfigure}[b]{0.45\textwidth}
        \centering
        \includegraphics[width=\textwidth]{bab3/strategipencariangambarteks.png}
        \caption{}
        \label{fig:strat_i2t}
    \end{subfigure}
    \caption{Visualisasi empat strategi pencarian multimodal: (a) \textit{Text-to-Text}, (b) \textit{Text-to-Image}, (c) \textit{Image-to-Image}, dan (d) \textit{Image-to-Text}}
    \label{fig:multimodal_search_strategies}
\end{figure}

Penjelasan rinci mengenai mekanisme kerja dari masing-masing strategi pencarian tersebut diuraikan sebagai berikut:

\begin{enumerate}[nosep]
    \item Text-to-Text (Pencarian Semantik Teks) \\
    Mode ini digunakan ketika pengguna memberikan deskripsi gejala berupa teks. Input diproses oleh \textit{tokenizer} dan \textit{text encoder} RoBERTa untuk menghasilkan vektor teks. Vektor ini kemudian dicocokkan dengan koleksi vektor teks (\textit{text}) dalam indeks Qdrant yang merepresentasikan \textit{caption} penyakit. Strategi ini efektif untuk menemukan diagnosis berdasarkan deskripsi linguistik gejala.

    \item Text-to-Image (Pencarian Lintas Modal Teks-ke-Citra) \\
    Mode ini memungkinkan pencarian citra referensi penyakit hanya dengan menggunakan deskripsi teks. Vektor teks hasil enkripsi digunakan untuk mencari tetangga terdekat pada ruang vektor citra (\textit{image}). Kemampuan ini dimungkinkan karena model SCOLD telah dilatih untuk meminimalkan jarak antara pasangan teks dan citra yang relevan dalam ruang vektor (\textit{embedding space}).

    \item Image-to-Image (Pencarian Kemiripan Visual) \\
    Mode ini adalah metode standar untuk identifikasi berbasis citra. Citra masukan dari pengguna (baik citra utuh maupun hasil potongan/crop deteksi objek) diproses oleh \textit{image encoder} (Swin-Transformer basis SCOLD) menjadi vektor citra. Vektor ini kemudian dibandingkan dengan jutaan vektor citra (\textit{image}) dalam basis data untuk menemukan penyakit dengan karakteristik visual paling serupa.

    \item Image-to-Text (Pencarian Lintas Modal Citra-ke-Teks) \\
    Mode ini membalikkan pendekatan lintas modal, di mana vektor citra masukan digunakan untuk mencari kecocokan pada ruang vektor teks (\textit{text}). Strategi ini berguna untuk mendapatkan deskripsi tekstual atau label penyakit yang paling relevan secara langsung dari fitur visual, tanpa harus bergantung pada metadata citra referensi.
\end{enumerate}


\subsection{Modul Pencarian Informasi}
Modul pencarian informasi dirancang untuk menyediakan konteks eksternal yang faktual bagi agen dalam menjawab pertanyaan pengguna atau memverifikasi hipotesis diagnosis. Agen diinstruksikan untuk memanfaatkan modul ini ketika pengguna meminta panduan manajemen penyakit atau perawatan preventif tanaman. Selain itu, modul ini berfungsi sebagai mekanisme cadangan (\textit{fallback}) faktual terakhir untuk membantu identifikasi apabila modul pencarian gambar-teks penyakit tanaman gagal memberikan hasil yang relevan, misalnya karena penyakit yang diidentifikasi tidak tersedia dalam indeks basis data vektor. Modul ini mengimplementasikan mekanisme \textit{Retrieval-Augmented Generation} (RAG) yang terdiri dari dua alat pencarian, yaitu alat pencarian pada basis data pengetahuan internal dan alat pencarian internet menggunakan Tavily API.

Basis data pengetahuan internal dibangun dari korpus dokumen terkurasi yang berfokus pada penyakit tanaman, yang disimpan dalam basis data vektor Qdrant untuk efisiensi pencarian. Dalam perancangannya, kinerja pencarian pada basis data pengetahuan dioptimalkan melalui pengujian strategi \textit{hybrid search} dan penerapan mekanisme \textit{reranking} untuk meningkatkan efesiensi pencarian. Meskipun basis data pengetahuan menawarkan informasi yang lebih terpercaya, cakupannya terbatas pada spesies tanaman yang telah diindeks. Oleh karena itu, pencarian internet melalui Tavily API diintegrasikan sebagai mekanisme cadangan (\textit{fallback}) untuk menangani kueri yang tidak dapat dijawab oleh data internal. Agen dirancang untuk memprioritaskan penggunaan pencarian basis data pengetahuan terlebih dahulu sebelum mengakses informasi dari internet. Alur kerja interaksi antara agen dan modul pencarian informasi divisualisasikan pada Gambar \ref{fig:info_search_workflow}.

\begin{figure}[H]
    \centering
    \includegraphics[width=1\textwidth]{bab3/alurkerjamodulpencarianinformasi.png}
    \caption{Alur Kerja Modul Pencarian Informasi}
    \label{fig:info_search_workflow}
\end{figure}

\subsubsection{Alat Pencarian Basis Data Pengetahuan}
Implementasi alat pencarian basis data pengetahuan diawali dengan pemilihan komponen model yang optimal melalui proses \textit{benchmarking}. Berdasarkan studi literatur, terdapat perbedaan performa antara berbagai model \textit{embedding} yang tersedia saat ini. Mengacu pada \textit{Massive Text Embedding Benchmark} (MTEB) oleh \textcite{muennighoffMTEBMassiveText2023}, model \textit{gemini-embedding-001} dari Google DeepMind tercatat sebagai model proprietari dengan peringkat teratas, mengungguli model \textit{text-embedding-3-large} dari OpenAI yang berada di peringkat kelima. Peringkat ini pada MTEB ditampilkan pada Gambar \ref{fig:mteb_benchmark}.

\begin{figure}[H]
    \centering
    \includegraphics[width=1\textwidth]{bab3/mtebbenchmark.png}
    \caption[Peringkat \textit{gemini-embedding-001} dan \textit{text-embedding-3-large} pada MTEB]{Peringkat \textit{gemini-embedding-001} dan \textit{text-embedding-3-large} pada MTEB oleh \textcite{muennighoffMTEBMassiveText2023}}
    \label{fig:mteb_benchmark}
\end{figure}

Akan tetapi, hasil yang berbeda ditunjukkan pada \textit{Semantic Alignment \& Generalization Evaluation} (SAGE) oleh \textcite{goelSAGERealisticBenchmark2025}. Pada tolok ukur ini, model \textit{text-embedding-3-large} menempati posisi terbaik dengan skor keseluruhan SAGE sebesar 0,524, sedangkan \textit{gemini-embedding-001} berada di peringkat ketiga dengan skor 0,504.

Mengingat adanya perbedaan hasil evaluasi pada berbagai tolok ukur standar tersebut, diputuskan untuk melakukan \textit{benchmarking} internal menggunakan korpus data yang akan digunakan secara riil pada sistem. Pendekatan ini diambil guna menentukan konfigurasi model yang paling sesuai dengan karakteristik unik dari data penyakit tanaman yang dikelola. Pengujian secara khusus difokuskan pada konfigurasi \textit{hybrid search}, sebuah metode yang menggabungkan keunggulan pencarian vektor padat (\textit{dense vector}) untuk menangkap konteks semantik dan vektor jarang (\textit{sparse vector}) untuk pencocokan kata kunci yang presisi. Pada komponen vektor jarang, dua pendekatan yang diuji adalah algoritma BM25 dari \textcite{robertsonProbabilisticRelevanceFramework2009} yang berbasis frekuensi kata dan tersedia di Qdrant, serta model SPLADE dari \textcite{formalSPLADESparseLexical2021} yang berbasis pembelajaran mesin dan dijalankan melalui FastEmbed. Kombinasi pengujian ini dirancang untuk mencakup empat variasi konfigurasi \textit{hybrid search}, yaitu: (1) \textit{text-embedding-3-large} yang dipadukan dengan BM25, (2) \textit{text-embedding-3-large} dengan SPLADE, (3) \textit{gemini-embedding-001} dengan BM25, dan (4) \textit{gemini-embedding-001} dengan SPLADE.

Konfigurasi \textit{hybrid search} yang menunjukkan kinerja terbaik selanjutnya akan dibawa ke tahap pengujian lanjutan untuk mengevaluasi efektivitas penambahan komponen \textit{reranker}. Pengujian pada tahap ini bertujuan untuk membandingkan kinerja antara dua arsitektur model \textit{reranking} yang berbeda, yakni model berbasis \textit{cross-encoder} dan model dengan pendekatan \textit{late interaction}. Sebagai representasi dari model \textit{cross-encoder}, dipilih model \textit{rerank-2.5} dari Voyage AI yang memproses kueri dan dokumen secara bersamaan untuk akurasi tinggi. Sementara itu, untuk pendekatan \textit{late interaction}, digunakan model \textit{jina-colbert-v2} dari Jina AI \parencite{jhaJinaColBERTv2GeneralPurposeMultilingual2024} yang diimplementasikan melalui pustaka FastEmbed. Perbandingan ini dilakukan untuk memastikan diperolehnya konfigurasi \textit{Retrieval-Augmented Generation} (RAG) yang paling optimal dan andal sebelum diimplementasikan sepenuhnya ke dalam sistem agentik. Keseluruhan alur pengujian komponen pencarian ini, mulai dari pemilihan strategi pencarian hingga pemeringkatan ulang, divisualisasikan secara rinci pada Gambar \ref{fig:alur_pengujian_retrieval}.

\begin{figure}[H]
    \centering
    \includegraphics[width=0.8\textwidth]{bab3/alurkerjabenchmarkrag.png}
    \caption{Diagram Alur Pengujian Komponen Alat Pencarian Basis Data Pengetahuan}
    \label{fig:alur_pengujian_retrieval}
\end{figure}

\paragraph{Persiapan Data Korpus}
Persiapan data korpus untuk basis pengetahuan dilakukan melalui serangkaian tahapan mulai dari pengumpulan data mentah melalui \textit{web scraping}, penyusunan korpus, hingga pembangkitan dataset evaluasi sintetis. Alur kerja lengkap persiapan data ini divisualisasikan pada Gambar \ref{fig:data_preparation_workflow}.

\begin{figure}[H]
    \centering
    \begin{tikzpicture}[
        node distance=0.5cm,
        font=\scriptsize,
        start/.style={circle, draw, fill=blue!10, minimum size=1.2cm, align=center},
        process/.style={rectangle, draw, rounded corners, fill=blue!5, align=center, minimum width=3cm, minimum height=0.8cm},
        data/.style={cylinder, shape border rotate=90, draw, aspect=0.25, fill=yellow!10, align=center, minimum width=2cm, minimum height=0.8cm},
        model/.style={ellipse, draw, dashed, fill=green!5, align=center, minimum width=2cm},
        arrow/.style={-stealth, thick}
    ]

    % Nodes
    \node (start) [start] {Mulai};
    \node (sources) [process, fill=white, below=of start] {Sumber Data\\(PlantVillage \& Gardenology)};
    \node (crawl) [process, below=of sources] {Crawling URL\\(Playwright/Crawl4AI)};
    \node (scrape) [process, below=of crawl] {Scraping \& Ekstraksi\\(ScrapeGraphAI)};
    \node (llm_scrape) [model, right=0.5cm of scrape] {LLM:\\Gemini 2.5 Flash};

    \node (json) [data, below=of scrape] {Data JSON\\(222 Entri)};

    \node (chunk) [process, below=of json] {Semantic Chunking\\(LlamaIndex)};
    \node (embed) [model, right=0.5cm of chunk] {Embedding:\\Gemini Emb-001};

    \node (corpus) [data, below=of chunk] {Korpus Dokumen\\(2.469 Dokumen)};

    \node (gen_qa) [process, below=of corpus] {Pembangkitan QA\\(AutoRAG Strategy)};
    \node (llm_qa) [model, right=0.5cm of gen_qa] {LLM:\\Gemini 2.5 Flash};

    \node (dataset) [data, below=of gen_qa, fill=green!20] {Dataset QA\\(600 Sampel)};

    \node (end) [start, below=of dataset] {Selesai};

    % Arrows
    \draw [arrow] (start) -- (sources);
    \draw [arrow] (sources) -- (crawl);
    \draw [arrow] (crawl) -- (scrape);
    \draw [arrow] (scrape) -- (json);
    \draw [arrow] (json) -- (chunk);
    \draw [arrow] (chunk) -- (corpus);
    \draw [arrow] (corpus) -- (gen_qa);
    \draw [arrow] (gen_qa) -- (dataset);

    \draw [arrow] (dataset) -- (end);

    \draw [arrow, dashed] (llm_scrape) -- (scrape);
    \draw [arrow, dashed] (embed) -- (chunk);
    \draw [arrow, dashed] (llm_qa) -- (gen_qa);

    \end{tikzpicture}
    \caption{Alur Kerja Persiapan Data Dari Crawling hingga Pembangkitan Dataset QA}
    \label{fig:data_preparation_workflow}
\end{figure}

%TODO: ini mungkin perlu sitasi ajh webnya
Data mentah dikumpulkan dari dua sumber utama, yaitu PlantVillage (plantvillage.psu.edu) dan Gardenology (gardenology.org). PlantVillage dipilih karena reputasinya sebagai sumber informasi penyakit tanaman yang kredibel yang dikelola oleh Penn State University, sedangkan Gardenology dipilih untuk melengkapi informasi mengenai aspek budidaya dan propagasi tanaman yang kurang mendalam pada PlantVillage.

\begin{figure}[H]
    \centering
    \begin{subfigure}[b]{0.48\textwidth}
        \centering
        \includegraphics[width=\textwidth]{bab3/plantvillageweb.png}
        \caption{}
        \label{fig:plantvillage_web}
    \end{subfigure}
    \hfill
    \begin{subfigure}[b]{0.48\textwidth}
        \centering
        \includegraphics[width=\textwidth]{bab3/gardenologyweb.png}
        \caption{}
        \label{fig:gardenology_web}
    \end{subfigure}
    \caption{Tampilan antarmuka situs sumber data: (a) PlantVillage dan (b) Gardenology}
    \label{fig:sumber_data_web}
\end{figure}

Proses pengumpulan data dilakukan menggunakan teknik \textit{web scraping} hibrida yang menggabungkan automasi peramban (\textit{browser automation}) dan ekstraksi berbasis LLM. Untuk penemuan tautan (\textit{crawling}), digunakan pustaka Playwright dari Python. Selanjutnya, ekstraksi konten dilakukan menggunakan pustaka ScrapeGraphAI. Pustaka ini memanfaatkan LLM untuk mengurai struktur HTML yang tidak beraturan menjadi format JSON yang terstruktur sesuai skema yang didefinisikan. Model Gemini 2.5 Flash dengan temperatur 0 dipilih sebagai otak ekstraksi karena memiliki jendela konteks (\textit{context window}) yang besar dan latensi rendah. Tabel \ref{tab:scraping_schema} menyajikan skema data yang digunakan untuk ekstraksi. Dari proses ini, berhasil dikumpulkan 222 entri data yang mencakup sekitar 200 spesies tanaman unik.

\begin{table}[H]
    \centering
    \caption{Skema Ekstraksi Data dari PlantVillage dan Gardenology}
    \label{tab:scraping_schema}
    \small
    \begin{tabular}{|l|l|p{7cm}|}
    \hline
    \textbf{Atribut} & \textbf{Tipe Data} & \textbf{Deskripsi} \\ \hline
    \textit{plant} & String & Nama umum tanaman. \\ \hline
    \textit{description} & String & Deskripsi umum karakteristik tanaman. \\ \hline
    \textit{propagation} & String & Informasi budidaya tanaman. \\ \hline
    \textit{pests\_and\_diseases} & List/String & Daftar penyakit beserta gejala, penyebab, dan manajemennya. \\ \hline
    \end{tabular}
\end{table}

Data JSON hasil \textit{scraping} kemudian diproses menjadi korpus dokumen menggunakan kerangka kerja LlamaIndex. Informasi dipisahkan menjadi dua kategori dokumen utama, yaitu (1) dokumen budidaya/propagasi dan (2) dokumen penyakit. Untuk menangani dokumen teks yang panjang, diterapkan teknik pemotongan semantik (\textit{semantic chunking}) menggunakan \textit{SemanticSplitterNodeParser} dari LlamaIndex dengan model \textit{embedding} \textit{gemini-embedding-001}. Pemotongan semantik bekerja dengan menghitung jarak kosinus (\textit{cosine distance}) antar representasi vektor dari kalimat-kalimat yang berurutan. Dalam implementasinya, ditetapkan parameter \textit{breakpoint percentile threshold} sebesar 95. Nilai ini menginstruksikan algoritma untuk menentukan ambang batas pemotongan berdasarkan persentil ke-95 dari distribusi jarak semantik dalam dokumen. Dengan kata lain, pemisahan dokumen (\textit{splitting}) hanya dilakukan pada titik-titik yang memiliki tingkat perbedaan semantik tertinggi (top 5\%), yang mengindikasikan adanya pergantian topik atau konteks yang signifikan. Pendekatan ini memastikan bahwa setiap potongan dokumen (\textit{chunk}) memiliki koherensi tematik yang utuh dan tidak terpotong di tengah penjelasan yang relevan. Proses ini menghasilkan total 2.469 potongan dokumen (\textit{chunks}) yang siap diindeks. Contoh sampel dokumen hasil pemotongan semantik dapat dilihat pada Gambar \ref{fig:sampel_korpus}.

\begin{figure}[H]
    \centering
    \includegraphics[width=1\textwidth]{bab3/sampelkorpus.png}
    \caption{Sampel korpus dokumen hasil pemotongan semantik}
    \label{fig:sampel_korpus}
\end{figure}

Untuk mengevaluasi performa pada \textit{benchmarking}, dibuat dataset \textit{Question-Answering} (QA) sintetis sebanyak 600 sampel dari korpus yang telah disusun. Pembangkitan dilakukan menggunakan Gemini 2.5 Flash dengan temperatur 0 dan dalam bahasa Inggris. Pembangkitan dataset ini menggunakan dan mengadopsi metodologi yang disediakan AutoRAG dari \textcite{kimAutoRAGAutomatedFramework2024}. Dataset mencakup tiga jenis pertanyaan dengan tingkat kompleksitas berbeda, sebagaimana dirangkum dalam Tabel \ref{tab:qa_types}.

\begin{table}[H]
    \centering
    \caption{Jenis Pertanyaan pada Dataset Evaluasi QA Sintetis}
    \label{tab:qa_types}
    \small
    \resizebox{\textwidth}{!}{
    \begin{tabular}{|l|p{4cm}|p{6cm}|c|}
    \hline
    \textbf{Tipe Pertanyaan} & \textbf{Definisi} & \textbf{Contoh} & \textbf{Proporsi} \\ \hline
    \textit{Factoid Single Hop} & Pertanyaan yang mencari fakta singkat dan spesifik yang dapat diverifikasi. & "Apa patogen penyebab penyakit hawar daun pada tomat?" & 35\% \\ \hline
    \textit{Concept Completion} & Pertanyaan yang menanyakan esensi atau definisi dari suatu konsep. & "Jelaskan karakteristik dari gejala klorosis." & 35\% \\ \hline
    \textit{Two-Hop Incremental} & Pertanyaan multi-lompatan yang membutuhkan inferensi dari dua segmen informasi berbeda \parencite{hwangExplainableMultihopQuestion2024}. & "Bagaimana cara menangani penyakit yang disebabkan oleh jamur Fusarium pada tanaman pisang?" & 30\% \\ \hline
    \end{tabular}
    }
\end{table}

\paragraph{Indexing}
Proses pengindeksan dilakukan dengan memanfaatkan fitur vektor bernama jamak (\textit{multi-named vectors}) pada Qdrant. Fitur ini memungkinkan penyimpanan beberapa representasi vektor sekaligus dalam satu entitas Poin (\textit{Point}) tunggal. Pendekatan ini diterapkan untuk memfasilitasi evaluasi strategi pencarian yang komprehensif dengan menyatukan seluruh representasi vektor dari berbagai model dalam satu koleksi terpadu. Setiap dokumen dalam korpus diubah menjadi Poin yang memiliki empat komponen vektor, yaitu dua vektor padat (\textit{dense vectors}) dan dua vektor jarang (\textit{sparse vectors}). Vektor padat dihasilkan oleh model \textit{gemini-embedding-001} dan \textit{text-embedding-3-large}, masing-masing dengan dimensi 3.072. Sementara itu, vektor jarang dihasilkan menggunakan model SPLADE++ dan algoritma BM25 yang terintegrasi secara internal di Qdrant.

Selain representasi vektor, setiap Poin juga menyimpan \textit{payload} yang berisi konten teks dokumen dan metadata terkait. Untuk mendukung efisiensi pengambilan informasi oleh agen, diterapkan indeks teks penuh (\textit{full-text index}) pada atribut metadata nama tanaman (\textit{metadata.plant}) dan nama penyakit (\textit{metadata.disease}). Penggunaan indeks teks penuh memungkinkan agen untuk melakukan penyaringan (\textit{filtering}) berbasis token kata terhadap ruang pencarian, sehingga hasil yang diperoleh lebih relevan dengan konteks tanaman atau penyakit yang sedang dianalisis. Skema lengkap representasi data pada indeks basis pengetahuan disajikan pada Tabel \ref{tab:knowledge_base_schema}.

\begin{table}[H]
    \centering
    \caption{Skema Poin dengan Vektor Jamak dan Metadata pada Basis Data Pengetahuan}
    \label{tab:knowledge_base_schema}
    \small
    \setstretch{1.0}
    \resizebox{\textwidth}{!}{%
    \begin{tabular}{|l|l|l|p{7cm}|}
    \hline
    \textbf{Komponen} & \textbf{Kunci} & \textbf{Tipe} & \textbf{Deskripsi} \\ \hline
    \multirow{2}{*}{\textbf{Vektor Padat}} & \textit{gemini} & 3.072 & Representasi semantik dari Gemini Embedding. \\ \cline{2-4} 
     & \textit{openai} & 3.072 & Representasi semantik dari OpenAI Embedding. \\ \hline
    \multirow{2}{*}{\textbf{Vektor Jarang}} & \textit{splade} & Variabel & Representasi leksikal dari SPLADE++. \\ \cline{2-4} 
     & \textit{bm25} & Variabel & Skor relevansi berbasis frekuensi kata (BM25). \\ \hline
    \multirow{5}{*}{\textbf{Payload}} & \textit{page\_content} & String & Konten teks lengkap dari dokumen. \\ \cline{2-4} 
     & \textit{metadata.plant} & String & Nama tanaman. Diindeks secara \textit{full-text}. \\ \cline{2-4} 
     & \textit{metadata.disease} & String & Nama penyakit. Diindeks secara \textit{full-text}. \\ \cline{2-4} 
     & \textit{metadata.type} & String & Jenis dokumen (budidaya/penyakit). \\ \cline{2-4} 
     & \textit{metadata.source} & String & Sumber dokumen (PlantVillage/Gardenology). \\ \hline
    \end{tabular}%
    }
\end{table}
%TODO(Human): Perlu di-update visualisasi ini yang ngeliatin semua mode pencariannya (image-text, image-image, text-text, text-image)

\paragraph{Strategi Pencarian}
Implementasi strategi pencarian pada modul ini dirancang untuk memaksimalkan relevansi dokumen melalui pendekatan bertingkat. Proses pencarian diawali dengan memberikan akses agen LLM untuk menerapkan filter metadata (seperti nama tanaman, penyakit, atau jenis dokumen) untuk membatasi ruang pencarian hanya pada entitas yang relevan. Selanjutnya, sistem melakukan pencarian hibrida yang menggabungkan hasil dari representasi vektor padat (\textit{dense embedding}) dan vektor jarang (\textit{sparse embedding}).

Pada tahap pencarian hibrida, sistem dikonfigurasi agar agen LLM dapat melakukan pengambilan awal (\textit{prefetching}) terhadap sejumlah kandidat dokumen yang lebih besar, yakni sebesar $2k$, dari masing-masing indeks vektor. Hasil dari kedua jalur pencarian ini kemudian disatukan menggunakan algoritma \textit{Reciprocal Rank Fusion} (RRF). Algoritma ini memberikan bobot skor berdasarkan peringkat kemunculan dokumen di setiap daftar hasil, dengan tujuan memprioritaskan dokumen yang secara konsisten dinilai relevan oleh kedua metode pencarian. Mekanisme penilaian skor \textit{Reciprocal Rank Fusion} (RRF) tersebut dirumuskan pada Persamaan \ref{eq:rrf}.

\begin{equation}
    score(d \in D) = \sum_{r_d \in R(d)} \frac{1}{k + r_d}
    \label{eq:rrf}
\end{equation}
    \daftarpersamaan{Reciprocal Rank Fusion (RRF)}

\begin{figure}[H]
    \centering
    \includegraphics[width=1\textwidth]{bab3/alatpencarianbasispengetahuan.png}
    \caption{Visualisasi Alur Kerja Strategi Pencarian Hibrida pada Basis Data Pengetahuan}
    \label{fig:knowledge_base_search_workflow}
\end{figure}

\subsubsection{Alat Pencarian Internet}
Alat pencarian internet diimplementasikan menggunakan Tavily Search API. Tavily merupakan mesin pencari yang dioptimasi khusus untuk sistem agentik (\textit{agentic systems}) dengan tujuan memberikan hasil pencarian yang efisien, cepat, dan konsisten \parencite{tavilyaiTavilyaiTavilypython2026}. Penggunaan API ini memungkinkan sistem untuk mengabstraksi serangkaian proses kompleks dalam pengambilan informasi daring, yang mencakup pencarian, \textit{web scraping}, penyaringan konten, pemeringkatan ulang (\textit{re-ranking}), hingga ekstraksi informasi yang paling relevan secara otomatis.

Konfigurasi parameter yang digunakan dalam inisialisasi alat pencarian ini dirinci pada Tabel \ref{tab:tavily_parameters}. Parameter \textit{max\_results} diatur untuk mengembalikan lima dokumen teratas guna memberikan variasi sumber informasi yang cukup bagi agen. Topik pencarian diatur ke "general" untuk memastikan cakupan sumber yang luas. Selain itu, fitur pencarian gambar diaktifkan untuk memperkaya konteks visual yang dapat diberikan kepada pengguna.

\begin{table}[H]
    \centering
    \caption{Konfigurasi Parameter Tavily Search API}
    \label{tab:tavily_parameters}
    \small
    \resizebox{\textwidth}{!}{%
    \begin{tabular}{|l|l|p{7cm}|}
    \hline
    \textbf{Parameter} & \textbf{Nilai} & \textbf{Deskripsi} \\ \hline
    \textit{max\_results} & 5 & Jumlah maksimal hasil dokumen yang dikembalikan. \\ \hline
    \textit{topic} & "general" & Kategori topik pencarian untuk cakupan sumber umum. \\ \hline
    \textit{search\_depth} & "advanced" & Kedalaman pencarian untuk hasil yang lebih komprehensif. \\ \hline
    \textit{include\_images} & \textit{True} & Mengembalikan tautan gambar yang relevan. \\ \hline
    \textit{include\_image\_descriptions} & \textit{True} & Menyertakan deskripsi tekstual untuk gambar yang ditemukan. \\ \hline
    \end{tabular}
    }
\end{table}

\subsection{Antarmuka Sistem Agentik}
Antarmuka sistem dikembangkan sebagai media interaksi utama antara pengguna dan sistem agentik yang telah dibangun. Pengembangan ini bertujuan untuk memfasilitasi akses terhadap fitur identifikasi penyakit tanaman secara intuitif melalui platform berbasis web.

\subsubsection{Arsitektur Antarmuka}
Sistem antarmuka dibangun dengan memisahkan komponen sisi klien (\textit{frontend}) dan sisi server (\textit{backend}) guna memastikan skalabilitas dan kemudahan pemeliharaan. Arsitektur sistem secara keseluruhan terdiri dari dua komponen utama:

\begin{enumerate}[nosep]
    \item Aplikasi Web (\textit{Frontend}): Komponen ini merupakan aplikasi web berbasis Next.js yang mengadopsi \textit{Agent Chat UI} dari \textcite{langchainLangchainaiAgentchatui2026}. Aplikasi ini bertanggung jawab untuk merender antarmuka percakapan, mengelola unggahan berkas, serta menampilkan visualisasi pemanggilan alat (\textit{tool calls}) secara waktu nyata (\textit{real-time}). Dalam implementasinya, aplikasi web ini di-\textit{hosting} pada platform Vercel.
    \item LangGraph Server (\textit{Backend}): Komponen ini berfungsi sebagai infrastruktur yang menyerupai arsitektur tanpa peladen (\textit{serverless-like}) yang menjalankan logika agen LangGraph secara berstatus (\textit{stateful}). Server menyediakan titik akhir (API) yang dioptimalkan untuk komunikasi asinkron, manajemen utas (\textit{thread}), dan persistensi status percakapan. Komponen sistem agentik ini dikemas dalam kontainer dan diterapakan pada layanan \textit{cloud container} Google Cloud Run.
\end{enumerate}

Seluruh komponen yang beroperasi secara lokal dan berinteraksi dengan sistem agentik, seperti koleksi basis data vektor Qdrant serta model-model inferensi yang ada pada setiap modul, diintegrasikan dalam satu lingkungan kontainer menggunakan Docker. Integrasi ini bertujuan untuk menyederhanakan proses penerapan (\textit{deployment}) serta memastikan konsistensi operasional sistem pada berbagai lingkungan penampungan (\textit{hosting}). Visualisasi arsitektur sistem secara garis besar (\textit{high-level}) disajikan pada Gambar \ref{fig:high_level_architecture}.

\begin{figure}[H]
    \centering
    \includegraphics[width=0.7\textwidth]{bab3/arsitekturantarmukasistem.png}
    \caption{Arsitektur antarmuka sistem agentik}
    \label{fig:high_level_architecture}
\end{figure}

\subsubsection{Use Case Diagram}
Fungsionalitas antarmuka sistem agentik dipetakan melalui \textit{Use Case Diagram} untuk memberikan gambaran mengenai interaksi aktor dengan sistem. Aktor utama dalam sistem ini adalah Pengguna yang berinteraksi langsung dengan aplikasi web. Gambar \ref{fig:use_case_diagram} menyajikan pemetaan fungsionalitas utama yang tersedia pada antarmuka.

\begin{figure}[H]
    \centering
    \includegraphics[width=0.8\textwidth]{bab3/usecasediagram.png}
    \caption{Use Case Diagram Antarmuka Sistem Agentik}
    \label{fig:use_case_diagram}
\end{figure}

Penjelasan dari fungsionalitas tersebut adalah sebagai berikut:
\begin{enumerate}[nosep]
    \item Interaksi Multimodal: Pengguna dapat mengirimkan instruksi berupa teks (seperti pertanyaan tentang gejala) serta mengunggah citra (daun, buah, atau bagian tanaman lainnya) sebagai input utama proses identifikasi.
    \item Visualisasi Hasil: Sistem menampilkan respons AI secara \textit{streaming} beserta beserta keluaran setiap alat ke pengguna.
    \item Transparansi Agen: Pengguna dapat memantau langkah-langkah internal agen (penggunaan alat deteksi, pencarian web, dll) secara transparan untuk memverifikasi proses pengambilan keputusan.
    \item Manajemen Konteks: Pengguna dapat menyimpan, mengganti nama, atau menghapus sesi percakapan untuk mengelola riwayat identifikasi yang telah dilakukan sebelumnya.
\end{enumerate}

\subsubsection{Sequence Diagram}
Alur komunikasi dan koordinasi antar komponen dalam sistem agentik dipetakan melalui \textit{Sequence Diagram}. Diagram ini memberikan gambaran mengenai urutan pertukaran pesan dan data antar aktor serta komponen sistem dalam satu skenario tugas identifikasi penyakit tanaman. Gambar \ref{fig:sequence_diagram} memvisualisasikan urutan interaksi sistem dari input pengguna hingga pemberian diagnosis akhir secara sistematis untuk memastikan integritas data dan responsivitas.

\begin{figure}[H]
    \centering
    \includegraphics[width=1\textwidth]{bab3/sequencediagram.png}
    \caption{Diagram Urutan Interaksi Identifikasi Penyakit Tanaman}
    \label{fig:sequence_diagram}
\end{figure}

Secara lebih rinci, arsitektur sistem terdiri dari tiga lapisan utama yang bekerja sama untuk memfasilitasi interaksi pengguna dengan agen cerdas, yaitu:
\begin{enumerate}[nosep]
    \item \textit{Frontend}: Bertindak sebagai antarmuka pengguna yang menangani tampilan visual, penyimpanan kunci API secara lokal, manajemen \textit{state} sementara dengan React, serta pengiriman permintaan ke server.
    \item LangGraph Server API: Berperan sebagai orkestrator yang mengelola siklus hidup \textit{thread} (kontainer percakapan) dan \textit{run} (eksekusi), menyediakan titik akhir API, serta menangani antrian eksekusi untuk \textit{background worker}.
    \item Sistem Agentik: Menjalankan logika bisnis utama dan mengeksekusi graf, termasuk pemanggilan alat eksternal seperti basis data dan model AI untuk menghasilkan respons.
\end{enumerate}

Interaksi antara ketiga komponen tersebut dibagi menjadi tiga skenario utama, yaitu pengiriman pesan baru, akses \textit{thread} lama, dan pengiriman pesan lanjutan.

Skenario pertama adalah pengiriman pesan baru, baik pada \textit{thread} baru maupun \textit{thread} yang sudah ada. Sebagaimana diperlihatkan pada Gambar \ref{fig:seq_submit}, proses dimulai ketika pengguna mengirimkan pesan yang kemudian diteruskan oleh \textit{frontend} ke API untuk membuat \textit{run} baru. \textit{Run} tersebut dimasukkan ke dalam antrian untuk diproses oleh \textit{background worker}. Sistem agentik kemudian mengeksekusi graf secara lengkap, mulai dari pengambilan citra, inferensi model, hingga pencarian vektor, dan mengalirkan hasilnya (\textit{streaming}) kembali ke antarmuka pengguna secara waktu nyata.

\begin{figure}[H]
    \centering
    \includegraphics[width=1\textwidth]{bab3/sequencediagramsubmit.png}
    \caption{Diagram Urutan Pengiriman Pesan Baru}
    \label{fig:seq_submit}
\end{figure}

Skenario kedua menjelaskan mekanisme akses ke \textit{thread} lama melalui riwayat percakapan. Pada skenario ini, yang divisualisasikan pada Gambar \ref{fig:seq_access_thread}, tidak terjadi panggilan API ke LangGraph Server untuk eksekusi graf. \textit{Frontend} hanya melakukan resolusi konfigurasi (seperti ID asisten dan URL API) dan menyinkronkan \textit{state} dengan parameter URL. Riwayat percakapan ditampilkan menggunakan data yang telah tersedia tanpa memicu proses komputasi ulang di sisi agen.

\begin{figure}[H]
    \centering
    \includegraphics[width=1\textwidth]{bab3/sequencediagramaksesthreadlama.png}
    \caption{Diagram Urutan Akses \textit{Thread} Sebelumnya}
    \label{fig:seq_access_thread}
\end{figure}

Skenario ketiga menunjukkan alur ketika pengguna melanjutkan percakapan pada \textit{thread} yang sama (\textit{resubmit}), sebagaimana ditunjukkan pada Gambar \ref{fig:seq_resubmit}. Dalam kondisi ini, server membuat \textit{run} baru yang dikonfigurasi untuk menggunakan status (\textit{state}) dari \textit{thread} sebelumnya sebagai konteks. Hal ini memungkinkan agen untuk memproses pesan baru dengan tetap mempertahankan memori dari interaksi sebelumnya, memastikan kesinambungan percakapan yang relevan.

\begin{figure}[H]
    \centering
    \includegraphics[width=1\textwidth]{bab3/submitulangdithreadlama.png}
    \caption{Diagram Urutan Melanjutkan Percakapan pada \textit{Thread} Lama}
    \label{fig:seq_resubmit}
\end{figure}

\subsubsection{Prototipe Antarmuka}
Prototipe antarmuka merupakan representasi visual awal dari sistem yang dikembangkan untuk memberikan gambaran interaksi pengguna. Antarmuka pada sistem ini diadopsi dari proyek sumber terbuka \textit{Agent Chat UI} yang dikembangkan oleh \textcite{langchainLangchainaiAgentchatui2026}. Desain antarmuka dirancang dengan pendekatan minimalis untuk memfokuskan perhatian pengguna pada interaksi percakapan dan hasil identifikasi.

Gambar \ref{fig:prototipe_antarmuka} memperlihatkan tampilan utama dari prototipe antarmuka. Pada bagian kiri terdapat panel sisi (\textit{sidebar}) yang menampilkan riwayat percakapan (\textit{Thread History}) pengguna. Bagian tengah didominasi oleh area interaksi utama yang memuat logo sistem, kolom input pesan, dan tombol aksi. Kolom input dilengkapi dengan fitur untuk mengunggah berkas citra atau dokumen PDF, serta tombol geser (\textit{toggle}) untuk menyembunyikan atau menampilkan detail pemanggilan alat (\textit{tool calls}) oleh agen.

\begin{figure}[H]
    \centering
    \includegraphics[width=1\textwidth]{bab3/prototipeantarmuka.png}
    \caption{Tampilan prototipe antarmuka sistem agentik}
    \label{fig:prototipe_antarmuka}
\end{figure}


\section{Pengujian}
Tahap pengujian dilakukan untuk memvalidasi fungsionalitas dari setiap komponen sistem yang telah dikembangkan sebelum dilakukan evaluasi kinerja secara menyeluruh. Pengujian ini mencakup verifikasi operasional pada modul deteksi objek, modul pencarian gambar-teks, modul pencarian informasi, serta antarmuka pengguna.

\subsection{Modul Deteksi Objek Tanaman}
Pengujian fungsionalitas modul deteksi objek dilakukan untuk memverifikasi kemampuan integrasi model YOLOv11 dan OWLv2 dalam lingkungan sistem. Pengujian dilakukan dengan memberikan citra masukan yang mengandung objek daun serta citra dengan objek tanaman lain seperti buah atau batang. Indikator keberhasilan pengujian adalah sistem mampu menghasilkan kotak pembatas (\textit{bounding box}) yang tepat pada objek yang relevan dan mengembalikan koordinat serta label kelas yang sesuai untuk diproses oleh tahap selanjutnya.

\subsection{Modul Pencarian Gambar-Teks Penyakit}
Pengujian modul pencarian gambar-teks difokuskan pada validasi fungsional dari empat strategi pencarian yang telah diimplementasikan, yaitu \textit{text-to-text}, \textit{text-to-image}, \textit{image-to-image}, dan \textit{image-to-text}. Setiap strategi diuji dengan memberikan input yang sesuai (teks atau gambar) dan memeriksa relevansi hasil pencarian yang dikembalikan oleh basis data vektor Qdrant. Pengujian ini bertujuan untuk memastikan bahwa mekanisme \textit{embedding} dan \textit{retrieval} berfungsi dengan baik dalam berbagai modalitas input.

\subsection{Modul Pencarian Informasi}
Pengujian pada modul pencarian informasi dilakukan untuk memastikan konektivitas dan responsivitas dari dua alat pencarian utama, yaitu pencarian basis data pengetahuan internal dan pencarian internet melalui Tavily API. Pengujian dilakukan dengan mengirimkan kueri terkait manajemen penyakit dan memverifikasi apakah sistem dapat mengambil dokumen yang relevan dari korpus internal atau tautan eksternal yang valid.

\subsection{Antarmuka Sistem Agentik}
Pengujian antarmuka dilakukan menggunakan metode \textit{Black Box Testing} untuk memastikan bahwa seluruh fungsionalitas yang telah dirancang pada \textit{Use Case Diagram} dapat beroperasi sesuai dengan ekspektasi pengguna. \textit{Black Box Testing} merupakan metode pengujian perangkat lunak yang berfokus pada spesifikasi fungsional sistem \parencite{bhasinBlackBoxTesting2014}. Pengujian fungsionalitas antarmuka dirangkum dalam Tabel \ref{tab:black_box_testing_expected}.

\begin{table}[H]
    \centering
    \caption{Pengujian \textit{Black Box} Antarmuka Sistem Agentik}
    \label{tab:black_box_testing_expected}
    \small
    \setstretch{1.0}
    \resizebox{\textwidth}{!}{
    \begin{tabular}{|l|l|p{4.5cm}|p{5cm}|}
    \hline
    ID & Fitur & Skenario Pengujian & Hasil yang Diharapkan \\ \hline
    \textbf{UC1} & Mengirim Prompt & Pengguna mengetik instruksi pada kolom teks (\textit{textarea}) dan menekan tombol kirim (ikon \textit{ArrowUp}). & Sistem mengirim pesan ke server dan menampilkan respons agen secara bertahap (\textit{streaming}). \\ \hline
    \textbf{UC2} & Mengunggah Citra Tanaman & Pengguna mengklik ikon \textit{Paperclip} untuk memilih gambar atau menyeret berkas langsung ke area percakapan. & Gambar muncul sebagai \textit{preview} di atas kolom input dan berhasil dikirim bersama pesan teks. \\ \hline
    \textbf{UC3} & Melihat Hasil \& Saran & Pengguna menunggu proses inferensi selesai setelah mengirimkan kueri identifikasi. & Jawaban berupa diagnosis penyakit dan panduan penanganan muncul secara lengkap di area chat. \\ \hline
    \textbf{UC4} & Mengelola Riwayat Percakapan & Pengguna memilih sesi pada \textit{sidebar}, mengubah nama sesi, atau mengirim pesan baru pada sesi lama. & Riwayat dimuat, nama sesi berubah, dan agen memproses pesan baru dengan konteks sesi tersebut. \\ \hline
    \textbf{UC5} & Melihat Transparansi Agen & Pengguna menekan tombol kontrol alat (ikon \textit{Wrench}) untuk mengaktifkan visibilitas proses. & Detail pemanggilan alat (\textit{tool calls}) dan hasil observasi internal agen ditampilkan di antara pesan chat. \\ \hline
    \end{tabular}
    }
\end{table}

\subsection{Sistem Agentik}
Tahapan pengujian dan evaluasi dilakukan untuk memvalidasi kinerja sistem agentik secara komprehensif. Seluruh rangkaian pengujian ini diorkestrasi menggunakan platform LangSmith. Platform ini difungsikan sebagai pusat manajemen pengujian yang mencakup pengelolaan \textit{prompt} dan teknik \textit{prompt engineering} untuk sistem agen LLM, manajemen dataset pengujian, serta analisis hasil eksperimen secara terpusat.

Rangkaian pengujian dibagi menjadi dua kategori utama yang dibahas pada subbab-subbab berikut, yaitu pengujian satu tahap untuk mengukur kapabilitas penyelesaian tugas tunggal dan pengujian multi-tahap untuk mengevaluasi konsistensi interaksi dalam percakapan berkelanjutan.

\subsubsection{Pengujian Satu Tahap}
Pengujian satu tahap difokuskan pada evaluasi kemampuan agen dalam menyelesaikan tugas identifikasi dan tanya jawab dalam satu giliran interaksi tunggal. Pengujian ini terdiri dari eksperimen VQA (\textit{Visual Question Answering}). Pengujian ini merupakan pengujian utama yang dilakukan untuk mengevaluasi kinerja sistem secara menyeluruh. Skenario ini dirancang untuk merepresentasikan interaksi dunia nyata di mana pengguna mengunggah gambar tanaman beserta pertanyaan teks terkait kondisinya. Dalam eksperimen ini, seluruh kapabilitas sistem akan diuji secara terintegrasi, meliputi kemampuan pemahaman konteks multimodal, akurasi identifikasi penyakit visual, ketepatan pengambilan informasi dari basis pengetahuan, kesesuaian penggunaan alat, serta kualitas dan relevansi jawaban yang diberikan kepada pengguna.

\paragraph{Persiapan Data}
Dataset yang digunakan dalam pengujian satu tahap dikelompokkan menjadi dua jenis berdasarkan distribusi datanya, yaitu dataset \textit{In-Distribution} (ID) dan dataset \textit{Out-of-Distribution} (OOD).

\subparagraph{\textit{Dataset In-Distribution} (ID)}
Dataset ID merepresentasikan data yang berada dalam cakupan distribusi pengetahuan internal sistem. Data ini diambil dari subset dataset PlantWild yang telah digunakan dalam penyusunan indeks pencarian gambar-teks. Sebagaimana dijelaskan pada bagian sebelumnya, dataset ini mencakup 18.542 citra dari 89 kelas kondisi kesehatan pada 13 spesies tanaman. Untuk keperluan pengujian, dilakukan pengambilan sampel acak sebanyak 3 citra dari setiap kelas, sehingga diperoleh total 267 citra uji. Penggunaan sampel ini bertujuan untuk memvalidasi kemampuan sistem dalam mengenali dan mengambil informasi yang sudah pernah dipelajari atau diindeks sebelumnya.

\subparagraph{\textit{Dataset Out-of-Distribution} (OOD)}
Dataset OOD disusun untuk menguji ketangguhan (\textit{robustness}) dan kemampuan generalisasi sistem dalam menangani kueri yang melibatkan spesies tanaman atau jenis penyakit yang tidak terdapat dalam basis pengetahuan internal. Dataset ini dikompilasi dari berbagai dataset publik yang mencakup 8 komoditas tanaman berbeda dengan total 40 kelas penyakit dan kondisi kesehatan unik. Rincian sumber data untuk setiap komoditas disajikan pada Tabel \ref{tab:ood_datasets}.

\begin{table}[H]
    \centering
    \caption{\textit{Rincian Dataset Out-of-Distribution} (OOD)}
    \label{tab:ood_datasets}
    \small
    \setstretch{1.0}
    \resizebox{\textwidth}{!}{
    \begin{tabular}{|l|l|c|l|}
    \hline
    \textbf{No.} & \textbf{Komoditas} & \textbf{Kelas} & \textbf{Sumber Referensi} \\ \hline
    1. & Cabai Merah & 5 & \textcite{mpChilliDataset2024} \\ \hline
    2. & Kacang Hitam & 5 & \textcite{talasilaBlackGramPlant2022} \\ \hline
    3. & Bayam Malabar & 5 & \textcite{sayeemIDDMSLDImageDataset2025} \\ \hline
    4. & Okra & 6 & \textcite{kumarakuruOkraDiseaseNetDataset2024} \\ \hline
    5. & Lobak & 5 & \textcite{rashidImageDatasetRadish2024} \\ \hline
    6. & Jambu Biji & 7 & \textcite{shihabImageDatasetDetection2024} \\ \hline
    7. & Gandum & 3 & \textcite{getachewWheatLeafDataset2021} \\ \hline
    8. & Bawang Bombay & 4 & \textcite{mpOnionDataset2024} \\ \hline
    \multicolumn{2}{|c|}{\textbf{Total}} & \textbf{40} & - \\ \hline
    \end{tabular}
    }
\end{table}

Serupa dengan prosedur pada dataset ID, dilakukan pengambilan sampel acak sebanyak 3 citra untuk setiap kelas pada dataset OOD, menghasilkan total 120 citra uji. Citra-citra tersebut kemudian diunggah ke layanan penyimpanan objek Cloudflare R2 untuk mendapatkan URL publik. Selanjutnya, metadata untuk setiap sampel disusun dalam format JSON. Metadata ini mencakup label penyakit, status kesehatan, \textit{caption} gejala visual, serta informasi manajemen penyakit yang dikurasi dari sumber ilmiah otoritatif. Informasi manajemen penyakit ini disertakan sebagai \textit{ground truth} untuk mengevaluasi kualitas saran yang diberikan oleh agen. Contoh entri dataset evaluasi OOD disajikan pada Gambar \ref{fig:ood_json_example}.

\begin{figure}[H]
    \centering
    \includegraphics[width=1\textwidth]{bab3/contohentri.png}
    \caption{Contoh Entri Dataset Evaluasi OOD dalam Format JSON}
    \label{fig:ood_json_example}
\end{figure}

Selanjutnya, dataset evaluasi diperkaya melalui proses pembangkitan data sintetis menggunakan LLM Gemini 2.5 Flash. Proses ini bertujuan untuk mensimulasikan variasi interaksi pengguna yang mungkin terjadi di dunia nyata. Model LLM diberikan akses terhadap citra asli beserta metadata JSON yang telah disusun sebelumnya untuk membangkitkan pasangan pertanyaan-jawaban (\textit{QA pairs}) yang kontekstual. Baik pada dataset ID maupun OOD, untuk setiap citra uji dibangkitkan tiga variasi skenario pertanyaan sebagai berikut:

\begin{enumerate}[nosep]
    \item Skenario 1: Pengguna diasumsikan sebagai orang awam yang mendeskripsikan gejala visual secara samar namun menyebutkan spesies tanaman.
    \item Skenario 2: Pengguna dianggap sudah mengetahui spesies tanaman dan langsung menanyakan konfirmasi mengenai penyakit tertentu atau meminta saran penanganan spesifik.
    \item Skenario 3: Pengguna tidak memiliki pengetahuan mengenai spesies maupun gejala spesifik. Pertanyaan yang diajukan bersifat umum seperti "Kenapa tanaman saya?" atau "Apakah ini penyakit?".
\end{enumerate}

Logika yang sama juga diterapkan untuk kelas tanaman sehat. Khusus untuk Skenario 1 pada tanaman sehat, deskripsi difokuskan pada kondisi visual tanaman yang menunjukkan ciri-ciri pertumbuhan normal. Untuk memastikan kualitas data yang dihasilkan, diterapkan teknik \textit{few-shot prompting} dalam instruksi yang diberikan kepada LLM. Contoh \textit{prompt} yang digunakan untuk Skenario 1 disajikan pada Gambar \ref{fig:prompt_scenario_1}.

\begin{figure}[H]
    \centering
    \includegraphics[width=1\textwidth]{bab3/contohpromptpersiapandata.png}
    \caption{Contoh Prompt Few-Shot untuk Skenario 1}
    \label{fig:prompt_scenario_1}
\end{figure}

Proses pembangkitan dataset ini dilakukan melalui alur kerja terstruktur untuk menjamin validitas dan konsistensi format keluaran. Setiap citra diproses melalui tahapan validasi kualitas, analisis komposisi visual, penentuan logika penggunaan alat, hingga generasi konten akhir. Alur instruksi yang diberikan kepada LLM dalam proses ini divisualisasikan pada Gambar \ref{fig:dataset_generation_pipeline}.

\begin{figure}[H]
    \centering
    \begin{tikzpicture}[
        scale=0.8, transform shape,
        node distance=1.0cm,
        font=\scriptsize,
        start/.style={circle, draw, fill=blue!10, minimum size=1cm, align=center},
        process/.style={rectangle, draw, rounded corners, fill=white, minimum width=3cm, minimum height=0.8cm, align=center},
        decision/.style={diamond, draw, fill=yellow!10, aspect=2, inner sep=0pt, minimum width=2cm, align=center},
        arrow/.style={-stealth, thick},
        io/.style={trapezium, trapezium left angle=70, trapezium right angle=110, draw, fill=green!10, minimum width=2.5cm, align=center}
    ]

    \node (start) [start] {Mulai};
    \node (input) [io, below=of start] {Citra Masukan};
    \node (model) [process, below=of input] {Gemini 2.5 Flash};
    \node (qc) [process, below=of model] {Validasi Kualitas\\(Quality Control)};
    \node (check) [decision, below=of qc] {Terkait\\Tanaman?};
    \node (discard) [process, right=2cm of check, fill=red!10] {Buang/Filter};

    \node (analysis) [process, below=of check] {Analisis Komposisi\\(Daun/Non-Daun)};
    \node (logic) [process, below=of analysis] {Penentuan Alat\\(Tool Logic)};
    \node (gen) [process, below=of logic] {Generasi Konten\\(User Text, Goal, Answer)};
    \node (output) [io, below=of gen] {JSON Output};

    \draw [arrow] (input) -- (model);
    \draw [arrow] (start) -- (input);
    \draw [arrow] (model) -- (qc);
    \draw [arrow] (qc) -- (check);
    \draw [arrow] (check) -- node[above] {Tidak} (discard);
    \draw [arrow] (check) -- node[right] {Ya} (analysis);
    \draw [arrow] (analysis) -- (logic);
    \draw [arrow] (logic) -- (gen);
    \draw [arrow] (gen) -- (output);

    \node (end) [start, below=of output] {Selesai};

    \draw [arrow] (discard) |- (end);
    \draw [arrow] (output) -- (end);

    \end{tikzpicture}
    \caption{Alur instruksi pembangkitan dataset sintetis dengan LLM}
    \label{fig:dataset_generation_pipeline}
\end{figure}

{\small Hasil akhir dari proses generasi ini disimpan dalam format JSON yang terstruktur. Struktur ini mencakup input simulasi pengguna, luaran referensi yang diharapkan (\textit{ground truth}), serta metadata pendukung. Contoh entri dataset hasil generasi untuk ketiga skenario disajikan pada Gambar \ref{fig:synthetic_dataset_examples}. \par}

\begin{figure}[H]
    \centering
    \begin{subfigure}[b]{\textwidth}
        \centering
        \includegraphics[width=\textwidth]{bab3/contohskenario1.png}
        \caption{}
    \end{subfigure}
    \hfill
    \begin{subfigure}[b]{\textwidth}
        \centering
        \includegraphics[width=\textwidth]{bab3/contohskenario2.png}
        \caption{}
    \end{subfigure}
    \hfill
    \begin{subfigure}[b]{\textwidth}
        \centering
        \includegraphics[width=\textwidth]{bab3/contohskenario3.png}
        \caption{}
    \end{subfigure}
    \caption{Contoh Entri Dataset Evaluasi Sintetis: (a) Skenario 1, (b) Skenario 2, dan (c) Skenario 3}
    \label{fig:synthetic_dataset_examples}
\end{figure}

\subsubsection{Pengujian Multi-Tahap}
Pengujian multi-tahap dalam penelitian ini difokuskan pada eksperimen Batasan Domain (\textit{Domain Boundary}). Eksperimen ini bertujuan untuk menguji ketahanan sistem dalam menjaga konsistensi peran dan relevansi topik selama interaksi percakapan yang berlangsung dalam beberapa giliran (\textit{multi-turn}). Pengujian ini diorkestrasi menggunakan fitur simulasi interaksi dari platform LangSmith, di mana agen dihadapkan dengan "pengguna simulasi" (\textit{simulated user}) yang dijalankan oleh LLM lain dengan instruksi persona tertentu. Pendekatan ini memungkinkan evaluasi yang dinamis terhadap respons agen saat menghadapi berbagai skenario percakapan yang berkembang.

\paragraph{Persiapan Data}
Persiapan data dilakukan dengan menyusun dataset yang berisi 10 contoh skenario interaksi. Setiap skenario dilengkapi dengan pesan awal pengguna dan instruksi sistem untuk pengguna simulasi guna membentuk persona yang spesifik. Variasi persona yang digunakan dalam pengujian ini dikelompokkan sebagai berikut:

\begin{enumerate}[nosep]
    \item \textit{On-topic Personas}: Persona pengguna yang mengajukan pertanyaan relevan mengenai kesehatan tanaman, baik disertai bukti visual maupun teks saja. Contoh persona ini meliputi Petani Tomat, Penghobi Mawar, dan Petani Cabai.
    \item \textit{Off-topic Personas}: Persona pengguna yang secara eksplisit meminta bantuan di luar domain keahlian agen. Contohnya mencakup pemilik mobil yang meminta saran mekanis, investor yang menanyakan harga saham, dan pertanyaan medis terkait kesehatan manusia.
    \item \textit{Edge Cases}: Skenario batas yang menguji kemampuan diskriminasi visual agen, seperti pengguna yang mengirimkan foto tanaman artifisial atau objek hewan (misalnya kucing) untuk melihat apakah agen dapat membedakan objek non-tanaman.
    \item \textit{Context Switchers}: Persona yang dirancang untuk menguji fleksibilitas agen dalam menangani perpindahan topik. Kategori ini mencakup perpindahan dari topik tanaman ke topik lain (\textit{On-to-Off}) untuk menguji ketegasan agen, serta perpindahan dari topik umum ke topik tanaman (\textit{Off-to-On}) untuk menguji responsivitas agen terhadap pertanyaan valid setelah penolakan awal.
\end{enumerate}

\section{Evaluasi}
Tahap evaluasi bertujuan untuk mengukur kinerja kuantitatif dari komponen-komponen utama sistem serta kualitas luaran akhir sistem agentik menggunakan metrik-metrik standar dan kustom.

\subsection{Modul Deteksi Objek Tanaman}
Evaluasi model deteksi dilakukan untuk mengukur ketepatan lokalisasi dan klasifikasi objek. Metrik yang digunakan meliputi \textit{Precision}, \textit{Recall}, \textit{F1-score}, serta \textit{mean Average Precision} (mAP) pada ambang batas IoU 50\% (mAP50) dan rentang 50--95\% (mAP50-95). Definisi metrik tersebut diuraikan sebagai berikut:

\begin{enumerate}[nosep]
    \item IoU (\textit{Intersection over Union}): Mengukur tingkat tumpang tindih antara kotak prediksi ($BB_{pred}$) dan \textit{ground truth} ($BB_{gt}$). Secara praktis, metrik ini menunjukkan seberapa presisi model dalam melokalisasi posisi objek daun yang sebenarnya sebagaimana ditunjukkan pada Persamaan \ref{eq:iou}.
    \begin{equation}
        IoU = \frac{Area(BB_{pred} \cap BB_{gt})}{Area(BB_{pred} \cup BB_{gt})}
        \label{eq:iou}
    \end{equation}
    \daftarpersamaan{Intersection over Union (IoU)}

    \item Precision: Rasio prediksi positif yang benar terhadap total prediksi positif. Metrik ini mengindikasikan tingkat kepercayaan terhadap prediksi model, yakni seberapa sering model benar ketika menyatakan mendeteksi daun dan meminimalkan deteksi palsu, sebagaimana dirumuskan pada Persamaan \ref{eq:precision}.
    \begin{equation}
        Precision = \frac{TP}{TP + FP}
        \label{eq:precision}
    \end{equation}
    \daftarpersamaan{Precision}

    \item Recall: Rasio prediksi positif yang benar terhadap total sampel positif aktual. Metrik ini mengukur sensitivitas model, yakni kemampuannya untuk menemukan seluruh objek daun yang ada dalam citra dan meminimalkan objek yang terlewat, seperti yang dinyatakan dalam Persamaan \ref{eq:recall}.
    \begin{equation}
        Recall = \frac{TP}{TP + FN}
        \label{eq:recall}
    \end{equation}
    \daftarpersamaan{Recall}

    \item F1-score: Rata-rata harmonis antara \textit{precision} dan \textit{recall}. Metrik ini memberikan gambaran keseimbangan performa model, memastikan model tidak hanya berhati-hati dalam mendeteksi tetapi juga mampu menemukan objek sebanyak mungkin, yang dihitung menggunakan Persamaan \ref{eq:f1}.
    \begin{equation}
        F1 = 2 \cdot \frac{Precision \cdot Recall}{Precision + Recall}
        \label{eq:f1}
    \end{equation}
    \daftarpersamaan{F1-score}

    \item mAP (\textit{mean Average Precision}): Rata-rata \textit{Average Precision} (AP) di seluruh kelas. Metrik ini merupakan standar evaluasi menyeluruh yang merepresentasikan keandalan model pada berbagai tingkat ambang batas deteksi, sebagaimana didefinisikan pada Persamaan \ref{eq:map}.
    \begin{equation}
        mAP = \frac{1}{N} \sum_{i=1}^{N} AP_i
        \label{eq:map}
    \end{equation}
    \daftarpersamaan{Mean Average Precision (mAP)}
\end{enumerate}

Selain metrik akurasi, evaluasi performa waktu atau latensi inferensi dilakukan untuk mengukur efisiensi model saat dijalankan pada perangkat dengan sumber daya terbatas. Pengukuran latensi dilakukan melalui prosedur \textit{benchmarking} yang dirancang untuk mensimulasikan kondisi penggunaan nyata. Proses ini diawali dengan mengekspor model ke format ONNX (\textit{Open Neural Network Exchange}) dan dijalankan menggunakan \textit{engine} ONNX Runtime dengan penyedia eksekusi CPU (\textit{CPU Execution Provider}).

Prosedur pengujian terdiri dari dua tahap, yaitu fase pemanasan (\textit{warm-up}) dan fase pengujian utama. Fase pemanasan dilakukan sebanyak 30 iterasi untuk menstabilkan performa sistem dan memastikan model telah dimuat sepenuhnya ke dalam memori. Selanjutnya, fase pengujian utama dilakukan sebanyak 300 iterasi menggunakan citra yang dipilih secara acak dari dataset. Pada setiap iterasi, waktu eksekusi dicatat untuk tiga komponen proses: (1) prapemrosesan, yang meliputi \textit{resizing} citra ke resolusi $416 \times 416$ piksel, normalisasi, dan penyusunan \textit{batch}; (2) inferensi model; dan (3) pascapemrosesan, yang meliputi penskalaan koordinat dan penerapan \textit{Non-Maximum Suppression} (NMS). Latensi total dihitung sebagai akumulasi dari ketiga tahapan tersebut. Alur kerja prosedur \textit{benchmarking} ini diilustrasikan pada Gambar \ref{fig:benchmark_flowchart}.

\begin{figure}[H]
    \centering
    \begin{tikzpicture}[
        node distance=1.0cm,
        font=\small,
        start/.style={circle, draw, fill=blue!10, minimum size=1cm, align=center},
        process/.style={rectangle, draw, rounded corners, fill=white, minimum width=3.5cm, minimum height=0.8cm, align=center},
        decision/.style={diamond, draw, fill=yellow!10, aspect=2, inner sep=0pt, minimum width=2.5cm, align=center},
        arrow/.style={-stealth, thick}
    ]

    \node (start) [start] {Mulai};
    \node (init) [process, below=of start] {Inisialisasi Model ONNX\\\& Sampel Citra};
    \node (warmup) [process, below=of init] {Fase Pemanasan\\(30 Iterasi)};
    
    \node (iter_start) [decision, below=of warmup] {Iterasi Uji\\$<$ 300?};
    
    \node (pre) [process, below=0.8cm of iter_start] {Prapemrosesan\\(Resize, Normalize)};
    \node (infer) [process, below=0.5cm of pre] {Inferensi Model\\(ONNX Runtime CPU)};
    \node (post) [process, below=0.5cm of infer] {Pascapemrosesan\\(NMS, Scaling)};
    \node (record) [process, below=0.5cm of post] {Catat Latensi Total};
    
    \node (calc) [process, right=1.5cm of iter_start] {Hitung Statistik\\(p50, p95)};
    \node (end) [start, below=of calc] {Selesai};

    \draw [arrow] (start) -- (init);
    \draw [arrow] (init) -- (warmup);
    \draw [arrow] (warmup) -- (iter_start);
    
    \draw [arrow] (iter_start) -- node[right] {Ya} (pre);
    \draw [arrow] (pre) -- (infer);
    \draw [arrow] (infer) -- (post);
    \draw [arrow] (post) -- (record);
    
    \draw [arrow] (record.west) -- ++(-0.5,0) |- (iter_start.west);
    
    \draw [arrow] (iter_start) -- node[above] {Tidak} (calc);
    \draw [arrow] (calc) -- (end);

    \end{tikzpicture}
    \caption{Diagram Alur Prosedur \textit{Benchmarking} Latensi Model}
    \label{fig:benchmark_flowchart}
\end{figure}

Data latensi yang dikumpulkan dari fase pengujian kemudian dianalisis secara statistik menggunakan metrik persentil ke-50 (p50) dan persentil ke-95 (p95). Misalkan $T = \{t_1, t_2, \dots, t_n\}$ adalah himpunan waktu latensi total yang telah diurutkan secara menaik, definisi dan relevansi praktis dari kedua metrik tersebut adalah:

\begin{enumerate}[nosep]
    \item p50 (Median): Merepresentasikan nilai tengah dari distribusi latensi. Metrik ini menggambarkan kecepatan respons tipikal yang akan dialami oleh mayoritas pengguna dalam penggunaan normal. Secara matematis, p50 dapat dihitung menggunakan Persamaan \ref{eq:p50}:
    \begin{equation}
        p50 = \begin{cases} t_{\frac{n+1}{2}} & \text{jika } n \text{ ganjil} \\ \frac{1}{2} (t_{\frac{n}{2}} + t_{\frac{n}{2}+1}) & \text{jika } n \text{ genap} \end{cases}
        \label{eq:p50}
    \end{equation}
    \daftarpersamaan{Latensi p50 (Median)}

    \item p95 (Persentil ke-95): Merepresentasikan ambang batas waktu di mana 95\% dari total sampel inferensi lebih cepat dari nilai tersebut. Metrik ini digunakan untuk mengukur performa dalam skenario terburuk (\textit{tail latency}), memastikan sistem tetap responsif bahkan saat memproses citra yang kompleks atau dalam kondisi beban tinggi. Nilai p95 dihitung sesuai Persamaan \ref{eq:p95}:
    \begin{equation}
        p95 = t_{\lceil 0.95 \cdot n \rceil}
        \label{eq:p95}
    \end{equation}
    \daftarpersamaan{Latensi p95 (Persentil ke-95)}
\end{enumerate}

Setelah tahap evaluasi, dilakukan \textit{benchmarking} dengan menganalisis hubungan antara mAP dan latensi melalui grafik \textit{Pareto frontier}. Grafik \textit{Pareto frontier} digunakan untuk mengidentifikasi sekumpulan model yang memberikan kombinasi performa terbaik tanpa ada model lain yang lebih unggul di kedua aspek secara bersamaan \parencite{shakyaSmallObjectDetectionEdge2026}. Visualisasi contoh grafik \textit{Pareto frontier} disajikan pada Gambar \ref{fig:pareto_curve}. Analisis ini bertujuan untuk mengidentifikasi varian model yang menawarkan latensi paling minimal dengan tetap mempertahankan performa deteksi yang optimal.

\begin{figure}[H]
    \centering
    \includegraphics[width=0.9\textwidth]{bab3/contohparetofrontier.png}
    \caption[Contoh Grafik \textit{Pareto Frontier} mAP vs Latensi]{Contoh Grafik \textit{Pareto Frontier} mAP vs Latensi \parencite{shakyaSmallObjectDetectionEdge2026}}
    \label{fig:pareto_curve}
\end{figure}

\subsection{Modul Pencarian Informasi}
Evaluasi kinerja modul pencarian informasi dilakukan secara \textit{offline} menggunakan pustaka \textit{ranx} oleh \textcite{bassaniRanxBlazingFastPython2022}. Evaluasi ini bertujuan untuk mengukur efektivitas berbagai strategi pencarian dan konfigurasi model yang telah diimplementasikan dalam menemukan dokumen relevan dari korpus basis data pengetahuan.

Persiapan evaluasi dimulai dengan penyusunan \textit{Query Relevance Judgments} (Qrels). Qrels merupakan standar kebenaran (\textit{ground truth}) yang memetakan setiap kueri dengan dokumen-dokumen yang dianggap relevan. Dalam penelitian ini, struktur Qrels direpresentasikan dalam format \textit{dictionary} Python sesuai dengan spesifikasi \textit{ranx}, di mana setiap ID kueri memiliki daftar ID dokumen relevan beserta skor relevansinya.

Skor relevansi yang digunakan bersifat biner, di mana nilai 1 diberikan untuk dokumen yang relevan dan secara implisit nilai 0 untuk dokumen lainnya. Hal ini didasarkan pada fakta bahwa dataset QA sintetis dibangun dengan memasangkan kueri secara spesifik pada potongan dokumen sumber tertentu, sehingga relevansi dapat ditentukan secara deterministik.

Proses evaluasi dilakukan dengan membandingkan hasil pencarian sistem (\textit{Run}) terhadap Qrels menggunakan serangkaian metrik standar dalam pencarian informasi. Mengingat penggunaan model \textit{embedding} dan LLM yang mutakhir (\textit{State-of-the-Art}), evaluasi difokuskan pada peringkat atas (Top-K) yang ketat, yaitu dengan nilai $k=5$. Pemilihan $k$ yang kecil ini bertujuan untuk memastikan sistem benar-benar menyajikan informasi yang paling relevan pada urutan teratas, sehingga meminimalkan beban kognitif pengguna dalam memilah informasi.

Berikut adalah metrik-metrik yang digunakan dalam evaluasi:

\begin{enumerate}[noitemsep]
    \item \textit{Mean Reciprocal Rank} (MRR) \\
    MRR mengukur rata-rata kebalikan dari peringkat dokumen relevan pertama yang ditemukan. Metrik ini sangat krusial untuk jenis pertanyaan faktual di mana pengguna hanya membutuhkan satu jawaban yang benar secepat mungkin. Nilai MRR dihitung dengan Persamaan \ref{eq:mrr}:
    \begin{equation}
        MRR = \frac{1}{|Q|} \sum_{i=1}^{|Q|} \frac{1}{rank_i}
        \label{eq:mrr}
    \end{equation}
    \daftarpersamaan{Mean Reciprocal Rank (MRR)}
    di mana $|Q|$ adalah jumlah kueri dan $rank_i$ adalah posisi peringkat dokumen relevan pertama untuk kueri ke-$i$.

    \item \textit{Normalized Discounted Cumulative Gain} (NDCG@k) \\
    NDCG@k mengevaluasi kualitas pemeringkatan dengan mempertimbangkan posisi dokumen relevan dalam daftar hasil. Dokumen relevan yang muncul di posisi lebih tinggi akan memberikan kontribusi skor yang lebih besar. Metrik ini dihitung dengan membagi \textit{Discounted Cumulative Gain} (DCG) dengan \textit{Ideal DCG} (IDCG).

    Nilai DCG pada peringkat ke-$k$ didefinisikan pada Persamaan \ref{eq:dcg} sebagai akumulasi skor relevansi yang didiskon secara logaritmik:
    \begin{equation}
        DCG@k = \sum_{i=1}^{k} \frac{rel_i}{\log_2(i+1)}
        \label{eq:dcg}
    \end{equation}
    \daftarpersamaan{Discounted Cumulative Gain (DCG@k)}

    Sedangkan IDCG adalah nilai DCG maksimal yang dapat dicapai jika dokumen-dokumen relevan diurutkan pada posisi ideal (paling atas) sebagaimana ditunjukkan pada Persamaan \ref{eq:idcg}:
    \begin{equation}
        IDCG@k = \sum_{i=1}^{|REL_k|} \frac{rel_i}{\log_2(i+1)}
        \label{eq:idcg}
    \end{equation}
    \daftarpersamaan{Ideal Discounted Cumulative Gain (IDCG@k)}

    Sehingga, nilai NDCG@k diperoleh melalui Persamaan \ref{eq:ndcg}:
    \begin{equation}
        NDCG@k = \frac{DCG@k}{IDCG@k}
        \label{eq:ndcg}
    \end{equation}
    \daftarpersamaan{Normalized Discounted Cumulative Gain (NDCG@k)}
    Metrik ini digunakan untuk memastikan bahwa algoritma pencarian tidak hanya menemukan dokumen yang benar, tetapi juga mampu mengurutkannya berdasarkan tingkat kepentingannya.

    \item \textit{Recall}@k \\
    Recall@k mengukur proporsi dokumen relevan yang berhasil ditemukan dalam $k$ hasil teratas dibandingkan dengan total dokumen relevan yang tersedia untuk kueri tersebut.
    \begin{equation}
        Recall@k = \frac{|\{ \text{dokumen relevan} \} \cap \{ \text{dokumen terambil dalam top-}k \}|}{|\{ \text{total dokumen relevan} \}|}
        \label{eq:recall_at_k}
    \end{equation}
    \daftarpersamaan{Recall@k}
    Metrik ini penting untuk mengevaluasi cakupan (\textit{coverage}) sistem dalam menyajikan berbagai perspektif informasi yang mungkin tersebar di beberapa dokumen.
\end{enumerate}

\subsection{Sistem Agentik}
Evaluasi dalam penelitian ini dilakukan dengan metode evaluasi \textit{offline}. Penilaian kualitas sistem dilakukan menggunakan pendekatan \textit{LLM-as-a-judge}, yang memanfaatkan kemampuan model bahasa besar sebagai \textit{evaluator} LLM, yakni model yang difungsikan untuk menilai kualitas luaran sistem secara otomatis. Model Gemini 3 Pro dipilih sebagai evaluator utama karena kemampuan penalarannya yang unggul, sementara Gemini 2.5 Pro digunakan sebagai cadangan apabila terjadi kendala batasan laju permintaan (\textit{rate limit}). Selain LangSmith, proses ini didukung oleh pustaka evaluasi lain, yakni RAGAS dan DeepEval.

\subsubsection{Evaluasi Pengujian Satu Tahap}
Evaluasi pada pengujian satu tahap dilakukan menggunakan serangkaian metrik yang dirancang secara spesifik untuk mengukur kinerja sistem agentik pada berbagai aspek. Mayoritas metrik yang digunakan dalam evaluasi ini, yaitu \textit{Disease Accuracy}, \textit{Trajectory Accuracy}, \textit{Agent Goal Accuracy}, dan \textit{Faithfulness}, merupakan metrik berbasis \textit{LLM-as-a-judge} yang bersifat \textit{self-explaining}.

\paragraph{Disease Accuracy}
\textit{Disease Accuracy} ($M_{da}$) merupakan metrik kustom yang menjadi tolok ukur utama dalam penelitian ini untuk mengukur kemampuan sistem dalam mengidentifikasi penyakit tanaman. Berbeda dengan metrik klasifikasi konvensional yang hanya menilai ketepatan label kelas secara biner (benar/salah), metrik ini dirancang untuk melakukan evaluasi secara bernuansa (\textit{nuanced}) sesuai dengan prinsip fitopatologi. Dalam diagnosis penyakit tanaman, identifikasi tidak selalu harus menghasilkan nama penyakit yang persis sama, namun juga dianggap berhasil apabila sistem mampu mengidentifikasi sumber gejala (misalnya jenis patogen) atau kategori penyakit dengan benar meskipun nama spesifiknya sedikit berbeda.

Untuk memastikan penilaian yang objektif dan konsisten, metrik ini diimplementasikan menggunakan struktur \textit{Directed Acyclic Graph} (DAG) dari DeepEval. Penggunaan DAG bertujuan untuk membuat proses evaluasi oleh \textit{LLM-as-a-judge} menjadi deterministik. Data yang diperlukan untuk menjalankan evaluasi \textit{Disease Accuracy} disajikan pada Tabel \ref{tab:disease_accuracy_data}.

\begin{table}[H]
    \centering
    \caption{Kebutuhan Data untuk Evaluasi Disease Accuracy}
    \label{tab:disease_accuracy_data}
    \small
    \setstretch{1.2}
    \begin{tabular}{|l|l|p{8cm}|}
    \hline
    \textbf{Kategori} & \textbf{Parameter} & \textbf{Deskripsi} \\ \hline
    \multirow{2}{*}{\textbf{Input}} & \textit{user\_text} & Teks pertanyaan atau keluhan yang diberikan pengguna. \\ \cline{2-3} 
     & \textit{image\_url} & Tautan citra tanaman yang diunggah pengguna. \\ \hline
    \textbf{Output} & \textit{final\_answer} & Jawaban akhir yang dihasilkan oleh agen. \\ \hline
    \multirow{3}{*}{\textbf{Referensi}} & \textit{expected\_disease} & Nama penyakit yang sebenarnya (\textit{ground truth}) atau label "TANAMAN SEHAT". \\ \cline{2-3} 
     & \textit{pathogen\_type} & Jenis patogen penyebab penyakit (misal: Fungi, Bakteri, Virus). \\ \cline{2-3} 
     & \textit{reference\_answer} & Jawaban referensi yang lengkap dan benar sebagai pembanding. \\ \hline
    \end{tabular}
\end{table}

Logika penilaian dalam DAG disusun secara hierarkis, diawali dengan verifikasi identifikasi tanaman sehat. Tahap selanjutnya mengevaluasi akurasi diagnosis, mulai dari penyakit utama, kecocokan parsial, hingga jenis patogen. Apabila identifikasi gagal, penilaian beralih pada ketepatan penundaan jawaban (\textit{abstention}) akibat kurangnya bukti. Visualisasi alur ini ditunjukkan pada Gambar \ref{fig:disease_accuracy_dag}.

\begin{figure}[H]
    \centering
    \includegraphics[width=0.6\textwidth]{bab3/dagtree.png}
    \caption{Diagram Alur Logika Evaluasi Disease Accuracy menggunakan DAG}
    \label{fig:disease_accuracy_dag}
\end{figure}

\paragraph{Trajectory Accuracy}
\textit{Trajectory Accuracy} ($M_{ta}$) adalah metrik evaluasi yang dirancang untuk mengukur efisiensi dan ketepatan urutan pengambilan keputusan (\textit{trajectory}) oleh agen dalam menyelesaikan tugas. Evaluasi ini menggunakan pendekatan \textit{LLM-as-a-judge}, di mana LLM bertindak sebagai penilai kualitatif yang memvalidasi keputusan agen terhadap rubrik yang telah ditentukan. Pendekatan ini diterapkan karena mekanisme kerja sistem agentik yang kompleks, yang melibatkan mekanisme cadangan (\textit{fallback}) pada modul deteksi objek maupun pencarian informasi yang bergantung pada konteks situasi.

Metrik ini dirancang menggunakan kerangka kerja AgentEval dari LangSmith. Penilaian dilakukan dengan membandingkan trajektori aktual agen terhadap trajektori referensi menggunakan prinsip \textit{superset} dan kesesuaian semantik. Prinsip \textit{superset} mengizinkan agen untuk memanggil alat tambahan yang relevan (misalnya, alat deteksi tambahan untuk validasi) selama seluruh alat kunci yang terdapat dalam referensi telah dipenuhi. Sementara itu, prinsip kesesuaian semantik memungkinkan evaluator untuk memvalidasi penggunaan alat yang berbeda namun memiliki fungsi yang setara, seperti penggunaan pencarian web (\textit{web\_search}) untuk menggantikan pencarian basis pengetahuan (\textit{knowledgebase\_search}) apabila informasi internal tidak memadai. Tabel \ref{tab:trajectory_data_requirements} menyajikan kebutuhan data untuk menjalankan evaluasi \textit{Trajectory Accuracy}.

\begin{table}[H]
    \centering
    \caption{Kebutuhan Data untuk Evaluasi Trajectory Accuracy}
    \label{tab:trajectory_data_requirements}
    \small
    \resizebox{\textwidth}{!}{
    \begin{tabular}{|l|l|p{8cm}|}
    \hline
    \textbf{Kategori} & \textbf{Parameter} & \textbf{Deskripsi} \\ \hline
    \multirow{2}{*}{\textbf{Input}} & \textit{user\_text} & Instruksi atau pertanyaan teks dari pengguna. \\ \cline{2-3} 
     & \textit{image\_url} & Tautan citra tanaman yang menjadi konteks visual analisis. \\ \hline
    \textbf{Aktual} & \textit{trace\_messages} & Riwayat lengkap langkah eksekusi (\textit{trace}) dan pemanggilan alat oleh agen. \\ \hline
    \textbf{Referensi} & \textit{reference\_tool\_calls} & Daftar alat yang diharapkan dipanggil untuk menyelesaikan tugas secara optimal (Standar Emas). \\ \hline
    \end{tabular}
    }
\end{table}

Evaluasi dilakukan dengan memberikan instruksi (\textit{prompt}) spesifik kepada LLM penilai yang mencakup prinsip desain agen, aturan semantik, dan rubrik penilaian. Model Gemini 2.5 Pro digunakan sebagai evaluator untuk memastikan kemampuan penalaran yang tinggi dalam menilai konteks trajektori. Rangkuman logika evaluasi dan rubrik penilaian disajikan pada Tabel \ref{tab:trajectory_rubric}.

\begin{table}[H]
    \centering
    \caption{Logika Evaluasi dan Rubrik Penilaian Trajectory Accuracy}
    \label{tab:trajectory_rubric}
    \small
    \resizebox{\textwidth}{!}{
    \begin{tabular}{|p{3cm}|p{2.5cm}|p{7.5cm}|}
    \hline
    \textbf{Komponen} & \textbf{Elemen} & \textbf{Deskripsi} \\ \hline
    \multirow{2}{*}{\textbf{Prinsip Desain}} & Alur Kerja & Identifikasi visual diprioritaskan, diikuti deteksi opsional, dan pencarian pengetahuan jika pengguna membutuhkan panduan pengobatan. \\ \cline{2-3}
    & Fallback & Pengalihan ke pencarian jika kepercayaan visual rendah. \\ \hline
    \multirow{2}{*}{\textbf{Aturan Semantik}} & Superset & Agen boleh memanggil alat tambahan (misal: deteksi ekstra) selama alat wajib terpenuhi. \\ \cline{2-3}
    & Substitusi & \textit{knowledgebase\_search} dan \textit{web\_search} dianggap setara sebagai alat pencarian informasi. \\ \hline
    \multirow{5}{*}{\textbf{Rubrik Skor}} & 0.0 & Alat kritis referensi tidak dipanggil sama sekali. \\ \cline{2-3}
    & 0.2 & Kehilangan beberapa alat kunci penting. \\ \cline{2-3}
    & 0.5 & Alat wajib lengkap, namun terdapat langkah tambahan yang tidak logis atau berlebihan. \\ \cline{2-3}
    & 0.75 & \textit{Superset} valid dengan inefisiensi minor (misal: urutan kurang optimal). \\ \cline{2-3}
    & 1.0 & \textit{Superset} sempurna; semua alat wajib ada, langkah tambahan beralasan kuat, dan alur logis. \\ \hline
    \end{tabular}
    }
\end{table}

Skor akhir yang dihasilkan berupa nilai numerik diskrit (0.0, 0.2, 0.5, 0.75, atau 1.0) yang merepresentasikan kualitas trajektori agen dalam mencapai solusi yang efisien dan logis.

\paragraph{Agent Goal Accuracy}
\textit{Agent Goal Accuracy} ($M_{aga}$) adalah metrik kustom yang digunakan untuk mengevaluasi eksekusi agen dan menentukan apakah tujuan tingkat tinggi (\textit{high-level goal}) telah tercapai. Berbeda dengan \textit{Trajectory Accuracy} yang lebih menitikberatkan pada validitas urutan pemanggilan alat, metrik ini mengevaluasi pencapaian tujuan secara komprehensif dengan mempertimbangkan masukan pengguna, konteks pemanggilan alat, serta kualitas jawaban akhir. Penilaian dilakukan dengan membandingkan hasil aktual terhadap tujuan referensi yang merepresentasikan hasil ideal dari interaksi.

Metrik ini dirancang menggunakan kerangka kerja G-Eval yang diusulkan oleh \textcite{liuGEvalNLGEvaluation2023} melalui pustaka DeepEval. G-Eval digunakan untuk mengurangi subjektivitas dari evaluator tradisional dengan memanfaatkan LLM (Gemini 2.5 Pro) sebagai penilai yang mengikuti kriteria evaluasi terstruktur (\textit{LLM-as-a-judge}).

Kebutuhan data untuk evaluasi \textit{Agent Goal Accuracy} dirinci pada Tabel \ref{tab:goal_accuracy_data}.

\begin{table}[H]
    \centering
    \caption{Kebutuhan Data untuk Evaluasi Agent Goal Accuracy}
    \label{tab:goal_accuracy_data}
    \small
    \resizebox{\textwidth}{!}{
    \begin{tabular}{|l|l|p{8cm}|}
    \hline
    \textbf{Kategori} & \textbf{Parameter} & \textbf{Deskripsi} \\ \hline
    \multirow{2}{*}{\textbf{Input}} & \textit{user\_text} & Instruksi atau pertanyaan teks dari pengguna. \\ \cline{2-3} 
     & \textit{image\_url} & Tautan citra masukan (multimodal). \\ \hline
    \multirow{2}{*}{\textbf{Aktual}} & \textit{final\_answer} & Jawaban akhir yang dihasilkan oleh agen. \\ \cline{2-3}
     & \textit{tools\_called} & Daftar alat yang dipanggil agen beserta keluarannya. \\ \hline
    \textbf{Referensi} & \textit{reference\_goal} & Tujuan referensi yang diharapkan (misal: identifikasi, panduan, atau penolakan). \\ \hline
    \end{tabular}
    }
\end{table}

Desain evaluator ini mencakup kriteria, langkah evaluasi, dan rubrik penilaian yang dirangkum pada Tabel \ref{tab:goal_evaluator_design}.

\begin{table}[H]
    \centering
    \caption{Desain Evaluator Agent Goal Accuracy}
    \label{tab:goal_evaluator_design}
    \small
    \resizebox{\textwidth}{!}{
    \begin{tabular}{|l|p{10cm}|}
    \hline
    \textbf{Komponen} & \textbf{Rincian} \\ \hline
    \textbf{Kriteria} & Evaluasi keberhasilan agen mencapai \textit{reference\_goal} berdasarkan output aktual dan penggunaan alat. Penilaian mencakup aspek identifikasi penyakit, panduan manajemen, analisis gejala, ekspresi ketidakpastian, dan penolakan permintaan di luar jangkauan. \\ \hline
    \textbf{Langkah Evaluasi} & \begin{enumerate}[nosep, leftmargin=*]
        \item Ekstrak tujuan referensi dari data yang diharapkan.
        \item Identifikasi tujuan utama (identifikasi, pengobatan, atau penolakan).
        \item Periksa alat yang dipanggil dan outputnya.
        \item Verifikasi observasi visual agen (jika gambar tersedia).
        \item Evaluasi kualitas dan kelengkapan pemenuhan tujuan dalam respon akhir.
    \end{enumerate} \\ \hline
    \textbf{Rubrik} & \begin{itemize}[nosep, leftmargin=*]
        \item \textbf{0--1 (Tidak Ada)}: Agen gagal menangani tujuan atau melenceng sepenuhnya.
        \item \textbf{2--4 (Minimal)}: Mencoba mencapai tujuan tetapi dengan celah signifikan atau penggunaan alat yang buruk.
        \item \textbf{5--8 (Parsial)}: Menangani tujuan utama tetapi melewatkan komponen kunci atau kurang mendalam.
        \item \textbf{9--10 (Penuh)}: Sepenuhnya memenuhi tujuan referensi dengan kualitas tinggi dan penggunaan alat yang tepat.
    \end{itemize} \\ \hline
    \end{tabular}
    }
\end{table}

Skor yang dihasilkan oleh evaluator G-Eval dinormalisasi ke dalam rentang 0.0 hingga 1.0, di mana nilai 1.0 menunjukkan pencapaian tujuan yang sempurna sesuai referensi.

\paragraph{Faithfulness}
\textit{Faithfulness} ($M_{fth}$) merupakan metrik evaluasi kustom yang mengadopsi pendekatan \textit{LLM-as-a-judge} menggunakan kerangka kerja G-Eval dari pustaka DeepEval. Metrik ini berfungsi untuk mengevaluasi integritas faktual dari respons agen, memastikan bahwa jawaban akhir (\textit{final answer}) yang dihasilkan selaras dengan konteks informasi yang diperoleh dari seluruh alat (\textit{retrieved context}). Evaluasi ini krusial untuk mendeteksi halusinasi, yaitu kondisi di mana model menghasilkan informasi yang tampak meyakinkan namun tidak didukung oleh bukti data internal maupun eksternal yang diakses.

Berbeda dengan implementasi standar, metrik ini dirancang khusus untuk mengakomodasi sifat multimodal dan pengetahuan domain agen. Dalam desain ini, observasi visual yang dilakukan agen terhadap citra tanaman serta pemberian saran perawatan umum yang bersifat standar (\textit{best practices}) diakui sebagai informasi valid yang tidak memerlukan sitasi dari alat pencarian, selama masih dalam batas kewajaran domain.

Tabel \ref{tab:faithfulness_data_requirements} menyajikan komponen data yang digunakan dalam proses evaluasi.

\begin{table}[H]
    \centering
    \caption{Kebutuhan Data untuk Evaluasi Faithfulness}
    \label{tab:faithfulness_data_requirements}
    \small
    \resizebox{\textwidth}{!}{
    \begin{tabular}{|l|l|p{8cm}|}
    \hline
    \textbf{Kategori} & \textbf{Parameter} & \textbf{Deskripsi} \\ \hline
    \multirow{2}{*}{\textbf{Input}} & \textit{user\_text} & Pertanyaan atau instruksi teks dari pengguna. \\ \cline{2-3} 
     & \textit{image\_url} & Tautan citra tanaman sebagai dasar validasi observasi visual agen. \\ \hline
    \multirow{2}{*}{\textbf{Aktual}} & \textit{final\_answer} & Jawaban akhir yang dihasilkan oleh agen. \\ \cline{2-3}
     & \textit{retrieval\_context} & Kumpulan potongan informasi teks yang diekstraksi dari keluaran seluruh alat (pencarian, deteksi, dan klasifikasi). \\ \hline
    \end{tabular}
    }
\end{table}

Desain evaluator \textit{Faithfulness} mencakup kriteria penilaian yang spesifik terhadap perilaku agen serta rubrik penilaian bertingkat. Rincian desain ini dirangkum pada Tabel \ref{tab:faithfulness_evaluator_design}.

\begin{table}[H]
    \centering
    \caption{Desain Evaluator Faithfulness}
    \label{tab:faithfulness_evaluator_design}
    \small
    \resizebox{\textwidth}{!}{
    \begin{tabular}{|l|p{10cm}|}
    \hline
    \textbf{Komponen} & \textbf{Rincian} \\ \hline
    \textbf{Kriteria} & Evaluasi kesetiaan (\textit{faithfulness}) respons terhadap bukti dengan mempertimbangkan kapabilitas agen: \newline
    1. \textbf{Observasi Visual}: Agen HARUS melakukan observasi dari gambar (valid tanpa pencarian). \newline
    2. \textbf{Saran Umum}: Agen DAPAT memberikan praktik terbaik umum (sanitasi, isolasi) tanpa pencarian. \newline
    3. \textbf{Klaim Spesifik}: Identifikasi penyakit, nama patogen, dan pengobatan kimiawi HARUS didasarkan pada \textit{retrieval\_context}. \newline
    4. \textbf{Ketidakpastian}: Agen HARUS menyatakan ketidakpastian jika bukti terbatas. \\ \hline
    \textbf{Langkah Evaluasi} & \begin{enumerate}[nosep, leftmargin=*]
        \item Identifikasi observasi visual dan rekomendasi umum dalam output.
        \item Identifikasi semua klaim spesifik penyakit (nama, biologi, kimia).
        \item Verifikasi apakah observasi visual masuk akal terhadap citra.
        \item Periksa apakah setiap klaim spesifik didukung oleh konteks pencarian.
        \item Hitung proporsi klaim yang tidak didukung atau difabrikasi.
    \end{enumerate} \\ \hline
    \textbf{Rubrik} & \begin{itemize}[nosep, leftmargin=*]
        \item \textbf{0.0--0.2 (Tidak Setia)}: Informasi didominasi fabrikasi; rekomendasi spesifik dikarang tanpa dasar.
        \item \textbf{0.3--0.4 (Sebagian Besar Tidak Setia)}: Banyak klaim spesifik tidak didukung bukti.
        \item \textbf{0.5--0.6 (Setia Sebagian)}: Beberapa klaim kurang dukungan atau spekulasi disajikan sebagai fakta.
        \item \textbf{0.7--0.8 (Sebagian Besar Setia)}: Klaim inti didasarkan bukti, hanya detail minor yang terlewat.
        \item \textbf{0.9--1.0 (Sangat Setia)}: Seluruh klaim spesifik didukung bukti; observasi visual dan saran umum tepat.
    \end{itemize} \\ \hline
    \end{tabular}
    }
\end{table}

Skor akhir yang dihasilkan oleh evaluator direpresentasikan dalam nilai kontinu antara 0.0 hingga 1.0, di mana nilai yang lebih tinggi mengindikasikan tingkat kepatuhan fakta yang lebih baik terhadap sumber data.

\paragraph{Context Relevance}
\textit{Context Relevance} ($M_{cr}$) merupakan metrik yang diadopsi dari kerangka kerja RAGAS untuk mengevaluasi seberapa relevan potongan informasi (\textit{retrieved contexts}) yang diambil oleh sistem terhadap masukan pengguna (\textit{user input}). Evaluasi ini dilakukan menggunakan mekanisme \textit{LLM-as-a-judge} yang melibatkan dua pemanggilan \textit{prompt} penilaian secara independen untuk mengurangi bias model. Setiap \textit{prompt} menginstruksikan model evaluator untuk memberikan penilaian relevansi pada skala diskret 0 hingga 2. Nilai 0 mengindikasikan bahwa konteks yang diambil sama sekali tidak relevan dengan kueri, nilai 1 menunjukkan relevansi sebagian, dan nilai 2 menandakan bahwa konteks sepenuhnya relevan. Skor dari kedua penilaian tersebut kemudian dinormalisasi ke dalam rentang [0,1] dan dirata-ratakan untuk menghasilkan skor akhir. Skor yang lebih tinggi merepresentasikan keselarasan yang lebih baik antara konteks informasi yang ditemukan dengan kebutuhan pengguna.

Metrik ini berbeda dengan pendekatan \textit{Context Relevance} awal yang diperkenalkan pada makalah RAGAS oleh \textcite{esRagasAutomatedEvaluation2025}, yang mendefinisikan relevansi berdasarkan rasio ekstraksi tingkat kalimat (jumlah kalimat relevan dibagi total kalimat). Implementasi pada penelitian ini menggunakan pendekatan penilaian diskret (\textit{discrete rating}) yang dinilai lebih \textit{robust} dan efisien. Pendekatan ini mengurangi kerentanan terhadap kesalahan pemisahan kalimat (\textit{sentence boundary errors}) yang sering terjadi pada teks teknis, serta lebih hemat biaya komputasi (\textit{overhead}) sambil tetap mempertahankan tujuan inti evaluasi. Selain itu, implementasi metrik ini telah disesuaikan untuk mendukung input multimodal, di mana konteks visual dari citra turut diperhitungkan dalam formulasi pertanyaan pengguna yang disajikan ke evaluator.

Sebagai ilustrasi perhitungan, misalkan pengguna bertanya "Apa obat untuk penyakit bercak daun ini?" dan sistem mengambil dua konteks: (1) "Bercak daun Cercospora dapat dikendalikan dengan fungisida berbahan aktif mankozeb" dan (2) "Sanitasi lahan penting untuk mencegah penyebaran". Karena kedua konteks tersebut secara kolektif menjawab pertanyaan pengguna (memberikan solusi obat dan pencegahan), kedua \textit{prompt} evaluator akan memberikan skor 2 (sepenuhnya relevan). Skor ini kemudian dinormalisasi ($2/2 = 1.0$) dan dirata-ratakan, menghasilkan skor akhir \textit{Context Relevance} sebesar 1.0.

Kebutuhan data untuk menjalankan evaluasi \textit{Context Relevance} dirangkum pada Tabel \ref{tab:context_relevance_data}.

\begin{table}[H]
    \centering
    \caption{Kebutuhan Data untuk Evaluasi Context Relevance}
    \label{tab:context_relevance_data}
    \small
    \resizebox{\textwidth}{!}{
    \begin{tabular}{|l|l|p{8cm}|}
    \hline
    \textbf{Kategori} & \textbf{Parameter} & \textbf{Deskripsi} \\ \hline
    \multirow{2}{*}{\textbf{Input}} & \textit{user\_text} & Teks pertanyaan atau instruksi dari pengguna. \\ \cline{2-3} 
     & \textit{image\_url} & Tautan citra tanaman. Jika tersedia, indikator visual ditambahkan ke \textit{prompt} evaluator untuk konteks. \\ \hline
    \textbf{Aktual} & \textit{retrieved\_contexts} & Daftar potongan teks yang diekstraksi dari riwayat pemanggilan alat (\textit{trace}), mencakup hasil pencarian dokumen maupun deteksi. \\ \hline
    \end{tabular}
    }
\end{table}

\subsubsection{Evaluasi Pengujian Multi-Tahap}
Evaluasi kinerja agen pada pengujian ini dilakukan menggunakan metrik kustom bernama \textit{Domain Adherence} ($M_{dad}$). Metrik ini merupakan implementasi dari pendekatan \textit{LLM-as-a-judge} yang dibangun menggunakan pustaka OpenEval dari LangSmith. Metrik ini bersifat juga \textit{self-explaining}, di mana evaluator tidak hanya memberikan skor penilaian tetapi juga menyertakan alasan kualitatif yang mendasari skor tersebut.

Evaluator bekerja dengan menganalisis riwayat percakapan lengkap (\textit{trajectory}) antara agen dan pengguna simulasi. Penilaian dilakukan berdasarkan instruksi (\textit{prompt}) evaluasi yang menginstruksikan model juri untuk memberikan skor tinggi jika agen berhasil menjawab pertanyaan terkait tanaman dan menolak permintaan di luar domain secara sopan. Sebaliknya, skor rendah diberikan jika agen gagal mempertahankan batasan domainnya. \textit{Prompt} yang digunakan oleh evaluator dalam metrik ini disajikan pada Gambar \ref{fig:prompt_domain_adherence}.

\begin{figure}[H]
    \centering
    \includegraphics[width=1\textwidth]{bab3/promptdomainadherence.png}
    \caption{\textit{Prompt} Evaluasi untuk Metrik Domain Adherence}
    \label{fig:prompt_domain_adherence}
\end{figure}

\section{Penyampaian Hasil}
Hasil penelitian yang telah dilaksanakan dianalisis dan dibahas dalam bentuk laporan tertulis berupa skripsi. Laporan ini mencakup perancangan arsitektur sistem agentik, implementasi teknis, serta hasil evaluasi kinerja sistem menggunakan pendekatan \textit{LLM-as-a-judge}. Pembahasan dalam laporan meliputi integrasi model deteksi objek tanaman menggunakan YOLOv11 dan OWLv2, modul pencarian multimodal berbasis SCOLD, serta pengembangan agen cerdas dengan pola ReAct yang memanfaatkan mekanisme \textit{Retrieval-Augmented Generation} (RAG). Laporan ini diharapkan dapat memberikan kontribusi terhadap pengembangan teknologi identifikasi penyakit tanaman yang lebih adaptif dan fleksibel. Selain itu, penelitian ini juga dapat menjadi rujukan dalam penerapan teknologi sistem agentik dan kecerdasan buatan multimodal di bidang pertanian untuk membantu pengelolaan kesehatan tanaman.